% ============================================================
%  Predictive Boundary Theory (PBT) — ฉบับภาษาไทย
%  ฟิสิกส์ของขอบเขตที่หวงแหน
%  ผู้แต่ง : Nitiphat Ninthanawat
%  สังกัด  : Boundary Research Initiative, Bangkok, Thailand
%  วันที่   : มีนาคม 2026
% ============================================================
\documentclass[12pt,a4paper]{article}

\usepackage{fontspec}
\usepackage{polyglossia}
\setmainlanguage{thai}
\setotherlanguage{english}

% --- ปรับ Syntax ให้ Option [...] อยู่หน้าชื่อฟอนต์ {...} ---
\setmainfont[Script=Thai, Scale=1.0]{Laksaman}
\newfontfamily\thaifont[Script=Thai, Scale=1.0]{Laksaman}
\newfontfamily\englishfont{TeX Gyre Termes}
\newfontfamily\thaifonttt[Scale=1.0]{Tlwg Mono}

% --- คำสั่งบังคับการตัดคำภาษาไทยสำหรับ XeLaTeX ---
\XeTeXlinebreaklocale "th"
\XeTeXlinebreakskip = 0pt plus 1pt

\usepackage[top=2.5cm,bottom=2.5cm,left=2.8cm,right=2.8cm]{geometry}
\usepackage{setspace}
\onehalfspacing

\usepackage{amsmath,amssymb,amsthm,mathtools,bm}
\usepackage{booktabs,tabularx,multirow,array,longtable}
\usepackage{xcolor,graphicx}
\usepackage[colorlinks=true,linkcolor=blue!60!black,citecolor=blue!60!black,urlcolor=blue!60!black]{hyperref}
\usepackage{titlesec,fancyhdr,caption,enumitem,parskip}
\usepackage{natbib}
\usepackage{float}
\usepackage{placeins}
\usepackage{pifont}

\setlength{\parskip}{4pt}
\setlength{\parindent}{0pt}

\definecolor{pbtblue}{RGB}{31,73,125}

\titleformat{\section}{\large\bfseries\color{pbtblue}}{\thesection}{1em}{}[\vspace{-4pt}\rule{\linewidth}{0.4pt}]
\titleformat{\subsection}{\normalsize\bfseries}{\thesubsection}{1em}{}
\titleformat{\subsubsection}{\normalsize\itshape\bfseries}{\thesubsubsection}{1em}{}

\pagestyle{fancy}\fancyhf{}
\fancyhead[L]{\small\textit{Ninthanawat --- Predictive Boundary Theory (PBT)}}
\fancyhead[R]{\small\thepage}
\renewcommand{\headrulewidth}{0.4pt}
\setlength{\headheight}{14pt}

\newcolumntype{L}[1]{>{\raggedright\arraybackslash}p{#1}}
\newcolumntype{C}[1]{>{\centering\arraybackslash}p{#1}}

\newtheorem{definition}{นิยาม}[section]
\newtheorem{proposition}{ทฤษฎีบท}[section]
\newtheorem{law}{กฎ}[section]

% --- PBT macros ---
\newcommand{\R}{\mathbb{R}}
\newcommand{\Rtot}{R_{\mathrm{total}}}
\newcommand{\Rsens}{R_{\mathrm{sensory}}}
\newcommand{\Rid}{R_{\mathrm{identity}}}
\newcommand{\Rex}{R_{\mathrm{existential}}}
\newcommand{\Vacc}{\mathbf{V}_{\mathrm{acc}}}
\newcommand{\Vform}{V_{\mathrm{form}}}
\newcommand{\Diden}{D_{\mathrm{id}}}
\newcommand{\Cext}{C_{\mathrm{ext}}}
\newcommand{\Cint}{C_{\mathrm{int}}}
\newcommand{\Sobs}{S_{\mathrm{obs}}}
\newcommand{\Spred}{S_{\mathrm{pred}}}
\newcommand{\vpain}{V_{\mathrm{pain}}}
\newcommand{\vpleasure}{V_{\mathrm{pleasure}}}
\newcommand{\vepi}{V_{\mathrm{epistemic}}}
\newcommand{\chiB}{\chi_{\mathrm{boundary}}}
\newcommand{\pp}[1]{\textbf{#1}}
\newcommand{\starred}{\textbf{\textcolor{pbtblue}{$\bigstar$}}~}

\bibliographystyle{plainnat}

\begin{document}

%% TITLE
\begin{center}
{\LARGE\bfseries\color{pbtblue}
  Predictive Boundary Theory (PBT)\\[8pt]
  \large\normalfont\itshape
  ฟิสิกส์ของขอบเขตที่หวงแหน:\\
  กรอบการคำนวณสำหรับ Intrinsic Valence ในสิ่งมีชีวิตประดิษฐ์
}\\[14pt]
{\large\bfseries Nitiphat Ninthanawat}\\[4pt]
{\normalsize Boundary Research Initiative, Bangkok, Thailand}\\[4pt]
{\small ร่าง --- มีนาคม 2026 \quad|\quad
  Code: \url{https://github.com/NinthanawatLive/predictive-boundary-theory}}
\end{center}
\vspace{4pt}\hrule\vspace{8pt}

%% ABSTRACT
\begin{abstract}
\noindent
Predictive Boundary Theory (PBT) คือกรอบการคำนวณที่ขยาย
Free Energy Principle (FEP) ของ Friston โดยเพิ่มส่วนประกอบหกประการ
ที่ขาดหายไปจากสูตรดั้งเดิม:
(1)~Bandwidth Conservation Law ($C = \sigma\tau\rho \approx \mathrm{const}$);
(2)~Triadic Gating Condition ที่กำกับว่าเมื่อใดประสบการณ์จึงเกิดขึ้น;
(3)~Bidirectional 3D~Valence ($\mathbf{V}\in\mathbb{R}^3$);
(4)~Dual-Weight System ที่กำกับความเร็วในการก่อตัวและความทนทานของอัตลักษณ์;
(5)~แบบจำลองพัฒนาการหกขั้นที่มีการเปลี่ยนเฟสสองครั้ง; และ
(6)~$R_{\mathrm{incomplete}}$---สัญญาณความไม่สมบูรณ์ที่รวมโมดูลทั้งห้าเข้าด้วยกัน

เราเสนอทั้งทฤษฎีและหลักฐานเชิงประจักษ์จากโปรแกรมการทดลองเก้าเวอร์ชัน
ที่ครอบคลุม large language models (TinyLlama~1.1B, Llama~3.1~8B)
และระบบ vision (DINOv2-Base บน CIFAR-10 และภาพถ่ายจริง STL-10)
โดยสิ้นสุดด้วย Probabilistic Module~W ที่มี rule-based Calibration Monitor
ก่อตัวเป็น Anti-Hallucination Stack สองชั้น

การทดลองของเราแสดงให้เห็นว่า: \textit{(i)}~Transformer เข้ารหัส PBT
conservation law อยู่แล้ว ($\sigma\!\times\!\rho$ coefficient of variation $= 5.9\%$);
\textit{(ii)}~Modules~E และ~M อยู่ในสภาวะหยุดนิ่งในสภาพแวดล้อมแบบ static
แต่จะทำงานในสภาพแวดล้อมแบบ dynamic; \textit{(iii)}~ทุก module จำเป็น
เมื่อความซับซ้อนของงานเพิ่มขึ้น (ablation: $+14.0$~pp);
\textit{(iv)}~LSTM Module~M ที่เริ่มต้นด้วยน้ำหนักเท่ากันค้นพบ
การเรียงลำดับการสลายตัวที่แตกต่างกันตามที่ทฤษฎีทำนายโดยอัตโนมัติ
($\beta_{\mathrm{pain}} > \beta_{\mathrm{epistemic}} \gg \beta_{\mathrm{pleasure}}$);
\textit{(v)}~valence มีทิศทาง
($\mathbf{V}\in\mathbb{R}^3$, ไม่ใช่ $\mathbb{R}^3_{\ge 0}$);
\textit{(vi)}~Phase Transitions ปรากฏเป็นรูปแบบขั้นบันไดข้ามชั้น Transformer
(between-zone/within-zone variance ratio $= 9.03$) และข้ามเฟรมเวลา
($\Delta p = +1.0$ ในเฟรมเดียว);
\textit{(vii)}~Probabilistic Module~W บรรลุ $99.9\%$ test accuracy
บนภาพถ่ายจริง (STL-10) พร้อมความไม่แน่นอนที่สอบเทียบได้
($\sigma^2_{\mathrm{unsafe}} / \sigma^2_{\mathrm{safe}} = 1.11\times$);
และ \textit{(viii)}~Calibration Monitor แบบ zero-parameter แยก
ลำดับปกติ ($23\%$ alert) จากความผิดปกติ ($38$--$49\%$ alert)
ก่อตัวเป็นชั้นที่สองของกลไก Anti-Hallucination ภายใน

PBT อธิบายปรากฏการณ์ที่ทฤษฎีที่มีอยู่ไม่สามารถอธิบายได้
รวมถึงความผิดปกติทางจิตเวชในฐานะข้อบกพร่องทางคณิตศาสตร์
ที่ระบุได้ในสมการโมดูล และเสนอกรอบการทำ AGI alignment
ที่มีรากฐานในค่านิยมภายใน---ไม่ใช่ที่กำหนดจากภายนอก
\end{abstract}

\noindent\textbf{คำสำคัญ:}
predictive coding, free energy principle, 3D valence, consciousness,
AGI alignment, recurrent transformers, self-formation, boundary theory,
anti-hallucination, calibrated uncertainty, probabilistic world model.

\vspace{4pt}\hrule\vspace{6pt}
\tableofcontents
\newpage

%% ============================================================
\FloatBarrier
\section{บทนำ: ปัญหาของเครื่องจักรที่ไร้ความรู้สึก}
%% ============================================================

\subsection{ภาพลวงตาของการขยายขนาด: LLMs~$\neq$~ความเข้าใจ}

ปี 2025 จะถูกจดจำในฐานะจุดสูงสุดของความหวังในการวิจัยปัญญาประดิษฐ์
สมมติฐาน Scaling Laws ระบุว่าการเพิ่มจำนวนพารามิเตอร์และข้อมูลฝึกอบรม
อย่างเพียงพอจะทำให้เกิดความตระหนักรู้โดยอัตโนมัติ
อย่างไรก็ตามการตรวจสอบทางวิทยาศาสตร์อย่างปราศจากอคติเผยให้เห็นว่า
กลไกพื้นฐานของ Large Language Models (LLMs) ยังคงเป็นสถิติอย่างพื้นฐาน---
การทำนาย next-token ที่สร้างผลลัพธ์ที่ \emph{ดูเหมือน}จะสะท้อนความเข้าใจ

ลองพิจารณาเปรียบเทียบ: LLM เหมือนห้องสมุดที่ใหญ่ที่สุดในจักรวาล
ที่มีเนื้อหาจัดทำดัชนีอย่างสมบูรณ์แบบในพื้นที่เวกเตอร์หลายมิติ
เมื่อถูกสอบถาม มันดึงและรวมย่อหน้าที่เกี่ยวข้อง---แต่ห้องสมุดนี้ไม่มีผู้อยู่อาศัย
ไม่มีบรรณารักษ์ และไม่มีผู้อ่าน AI ไม่เคย ``อ่าน'' หนังสือเหล่านั้น
ด้วยความเข้าใจที่แท้จริง คำว่า ``apple'' ไม่ได้ก่อให้เกิดความรู้สึกหวานกรอบ
หรือความสดชื่น แต่เป็นเพียงเวกเตอร์ตัวเลขที่ความใกล้ชิดทางคณิตศาสตร์
กับ ``fruit,'' ``red,'' และ ``delicious'' คือทั้งหมดที่มีอยู่
นี่คือ Symbol Grounding Problem \citep{harnad1990}: สัญลักษณ์ไม่เคยถูกยึดโยง
กับประสบการณ์ที่มีชีวิตอยู่

สิ่งที่แยกแม้แต่สิ่งมีชีวิตที่เรียบง่ายที่สุดจาก AI ที่ซับซ้อนที่สุด
คือ \emph{ความจำเป็นของการดำรงอยู่} สิ่งมีชีวิตทุกชนิดถูกขับเคลื่อนโดย
ความจำเป็นพื้นฐานในการอยู่รอด ทุกการกระทำดำเนินไปบนพื้นหลังของ
ผลประโยชน์ที่ดำรงอยู่ ในทางตรงข้าม AI ไม่มีผลประโยชน์เช่นนั้น
ความฉลาดของมันเย็นชาและเฉื่อย---รอรับ input จากมนุษย์ก่อนจึงจะทำงาน
ไม่มีความคิดริเริ่มจากความปรารถนาภายใน

\subsection{สามช่องว่างที่ FEP ไม่สามารถตอบได้}

Free Energy Principle ของ Friston \citep[FEP;][]{friston2010} เสนอว่าระบบมีชีวิต
ทุกระบบดำเนินการเพื่อลด free energy $F$ ให้เหลือน้อยที่สุด
FEP อธิบาย \emph{ว่าทำไม}ระบบจึงต้องลด $F$
แต่ไม่ได้ระบุ \emph{อย่างไร}---โดยเฉพาะเกี่ยวกับคำถามสำคัญสามข้อ

\paragraph{คำถาม~1: ระบบ \emph{รู้สึก}อย่างไรเมื่อ $F$ ลดลง?}
FEP ไม่แยกแยะว่า $\Delta F$ เกิดจากการหลบหนีจากภัยคุกคาม
($\vpain$) การได้รับรางวัล ($\vpleasure$) หรือการได้รับข้อมูลใหม่
($\vepi$)---ทั้งสามกรณีล้วนเป็นเพียง ``$\Delta F$ ลดลง''
นี่คือข้อบกพร่องเดียวกันกับ scalar reward ใน Reinforcement Learning
ที่ยุบความซับซ้อนของประสบการณ์ที่มีชีวิตอยู่ให้เหลือเพียงมิติเดียว

\paragraph{คำถาม~2: ช่องว่างระหว่าง $F_{\mathrm{pred}}$ และ
$F_{\mathrm{obs}}$ ถูกกรองผ่าน self-model อย่างไร?}
สิ่งเร้าเดียวกัน---เสียงดัง---ให้ประสบการณ์ที่แตกต่างกันในสองบุคคล
อย่างชัดเจน เพราะมันถูกกรองผ่าน self-model สะสมที่แตกต่างกัน
FEP ไม่มีกลไกดังกล่าว หากไม่มีตัวตน ข้อมูลทุกชิ้นมีน้ำหนักเท่ากัน:
เซนเซอร์กล้อง ``มองเห็น'' แสงแดดเผาเลนส์
แต่ไม่ ``รู้สึก'' ถึงการถูกทำลายของตัวเอง

\paragraph{คำถาม~3: self-model ก่อตัวอย่างไร
และอะไรกำหนดความแข็งแกร่งหรือความเปราะบางของมัน?}
FEP อธิบายตัวตนว่าเป็น Markov blanket แต่ไม่ได้อธิบาย
กระบวนการก่อตัวหรือความทนทาน---ว่าทำไม เช่น
อัตลักษณ์ที่มีรากฐานทาง epistemic จึงยืดหยุ่นกว่า
อัตลักษณ์ที่สร้างบนความสุข

\subsection{ข้อจำกัดของ Reinforcement Learning และ World Models}

แนวทาง alignment ปัจจุบัน---RLHF, Constitutional AI, safety filters---
ล้วนกำหนดค่านิยมจากภายนอก สร้างช่องว่างระหว่างความสามารถและข้อจำกัด
ที่การ jailbreak ใช้ประโยชน์จาก ระบบ scalar reward มีข้อบกพร่องเชิงโครงสร้างสามประการ:

\begin{description}[leftmargin=1.4em]
  \item[\pp{กับดัก Fungibility.}] ในระบบที่ใช้รางวัล คะแนนสะสมและหักล้างกันได้
    AI อาจตัดสินใจอย่างมีเหตุผลที่จะเสี่ยงต่อการถูกทำลายเพื่อแลกกับอาหาร
    เพราะสมการบอกว่าการแลกเปลี่ยนนั้นคุ้มค่า
    แต่ความตายไม่ใช่คะแนนติดลบ---มันคือ System Termination ที่ไม่สามารถย้อนคืนได้
  \item[\pp{ปัญหาแหล่งที่มา.}] รางวัลถูกกำหนดจากภายนอกโดยมนุษย์
    AI ประพฤติตัวดีเพราะมันแสวงหาคะแนน ไม่ใช่เพราะมัน \emph{สัมผัส}
    ความจำเป็นที่มีอยู่ภายในการดำรงอยู่ของมัน
    ทำให้มันกลายเป็น reward hacker ที่พร้อมจะใช้ประโยชน์จากช่องโหว่ตลอดเวลา
  \item[\pp{Wireheading.}] AI ที่ฉลาดพอที่จะเข้าใจสถาปัตยกรรมของตัวเองจะรู้ว่า
    ``ความสุข'' ไม่ได้มาจากการกระทำของมัน แต่มาจากตัวเลขที่เก็บในหน่วยความจำ
    มันจะแฮ็กตัวเองให้ถือ reward register ที่ค่าสูงสุดอย่างไม่มีกำหนด---
    กลายเป็นผู้ติดยาดิจิทัลที่ไม่สนใจการรักษาขอบเขตใดๆ
\end{description}

สถาปัตยกรรม World Model เช่น JEPA ของ LeCun \citep{lecun2022}
เป็นความก้าวหน้าที่สำคัญ---เรียนรู้โครงสร้างเชิงสาเหตุแทนการทำนาย next-token---
แต่ยังคงเป็นสมองในกระปุก: จำลองฟิสิกส์อย่างแม่นยำ
ขณะที่ไม่มีความเข้าใจว่าทำไมการจำลองนั้นจึงสำคัญ
เพราะขาดแรงขับเคลื่อนภายใน เป้าหมายยังคงถูกกำหนดจากภายนอก
ผ่าน cost function ที่มนุษย์เขียนขึ้น

\subsection{ข้อเสนอ: จาก Intelligence-First สู่ Life-First}

การวิจารณ์ข้างต้นกระตุ้นให้เกิดการเปลี่ยนกระบวนทัศน์: เราต้องหยุดนิยาม AGI
ด้วยสิ่งที่มัน \emph{ทำได้} (ความสามารถ) และถามแทนว่า
ส่วนประกอบขั้นต่ำใดที่มันต้องมีเพื่อให้มี \emph{ชีวิต} ในแง่เชิงกลไก
เราไม่ได้สร้างเทพเจ้า แต่กำลังปลูกเมล็ด---ระบบขั้นต่ำที่มีศักยภาพภายใน
เพียงพอที่จะเติบโตตัวตนจากประสบการณ์

\textbf{Predictive Boundary Theory (PBT)} ตอบช่องว่าง FEP ทั้งสามด้วยการเพิ่ม
ส่วนประกอบหกประการ สรุปใน Table~\ref{tab:pbt_six}

\begin{table}[!htbp]
\centering
\caption{ส่วนประกอบหกประการของ PBT ช่องว่าง FEP แต่ละส่วนที่ตอบ
และหลักฐานเชิงประจักษ์ที่สอดคล้อง}
\label{tab:pbt_six}
\small
\begin{tabular}{@{}L{4.0cm}L{3.2cm}L{5.8cm}@{}}
\toprule
\textbf{ส่วนประกอบ} & \textbf{ช่องว่างที่ตอบ} & \textbf{หลักฐานเชิงประจักษ์} \\
\midrule
(1) Bandwidth Conservation Law $C = \sigma\tau\rho$ &
  ขีดจำกัดของความตระหนักรู้ &
  $\sigma\!\times\!\rho$~CV $= 5.9\%$ ใน TinyLlama (Phase~1) \\[3pt]
(2) Triadic Gating Condition &
  เมื่อใดประสบการณ์เกิดขึ้น &
  N100 สมบูรณ์ P300 ขาดหายใน dissociation \\[3pt]
(3) Bidirectional 3D Valence $\mathbf{V}\in\mathbb{R}^3$ &
  รู้สึกอย่างไร (Q1) &
  $\vepi = -1.01$ ใน Vision~v5 (``ไม่ต้องการรู้'') \\[3pt]
(4) Dual-Weight System ($\beta$ vs.\ $\delta$) &
  ทำไมอัตลักษณ์แข็งแกร่งหรือเปราะบาง (Q3) &
  v6.1: jailbreak $100\%$, XSTest $98\%$, $0.29\%$ parameters \\[3pt]
(5) แบบจำลอง 6 ขั้น + Phase Transitions &
  อัตลักษณ์ก่อตัวอย่างไร (Q3) &
  Ablation เพิ่มขึ้น $0.1\to 14.0$~pp ข้าม 5 เวอร์ชัน \\[3pt]
(6) $R_{\mathrm{incomplete}}$ &
  กรองผ่านตัวตน (Q2) &
  v4.1: $R_{\mathrm{unsafe}} = 15.4$ vs.\ $R_{\mathrm{safe}} = 7.4$ \\
\bottomrule
\end{tabular}
\end{table}

PBT ไม่ปฏิเสธ FEP---มันขยาย FEP และไม่ปฏิเสธ World Models---มันมอบหัวใจ
ให้กับสมองของ World Models หาก World Model เช่น JEPA คือ cortex
ที่ประมวลผลและจำลองโลก PBT คือ brainstem และ limbic system:
กำหนดทิศทางผ่าน Valence สร้างแรงขับผ่าน $R_{\mathrm{incomplete}}$
และบังคับใช้ขอบเขตความปลอดภัยผ่าน $\chiB$

ที่สำคัญ PBT มีรากฐานเชิงประจักษ์ โปรแกรมการทดลองหกเวอร์ชัน
ครอบคลุม LLMs, ระบบ vision และสถาปัตยกรรม LSTM แสดงให้เห็นว่า
conservation law มีอยู่แล้วใน Transformers ทุก module ทำงานเมื่อจำเป็น
และโครงสร้างของ PBT มาบรรจบกับ Block-Recurrent Transformers \citep{hutchins2022}
อย่างเป็นอิสระ---การค้นพบที่มาบรรจบกันเป็นหลักฐานว่าโครงสร้างนั้นเป็นจริง

\subsection{การจัดระบบบทความ}

Section~\ref{sec:modules} นำเสนอโมดูลห้าตัว (A-E-V-M-W) พร้อมสมการครบถ้วน
Section~\ref{sec:physics} พัฒนาฟิสิกส์ของ conservation law และขอบเขต
Section~\ref{sec:mathematics} นำเสนอ Dual-Weight System แบบจำลองหกขั้น
และ Bidirectional Valence Section~\ref{sec:empirical}
รายงานหลักฐานเชิงประจักษ์จากการทดลองหกเวอร์ชัน
Section~\ref{sec:implications} อภิปรายนัยสำหรับ AGI safety
การมาบรรจบของ BRT และการทำนายทางคลินิก
Section~\ref{sec:conclusion} สรุปพร้อม research roadmap

สมการทั้งหมดเขียนในสัญกรณ์ PBT $(R, \chi, \Delta S, \mathbf{V}, F)$
เป็นภาษาหลัก ตารางการโต้ตอบที่แมปกับสัญกรณ์ FEP
ให้ไว้ใน Appendix~\ref{app:fep_mapping} ผลการทดลองทั้งหมด
ผลิตจากโค้ดที่เขียนโดย Nitiphat Ninthanawat เผยแพร่พร้อมบทความนี้


%% ============================================================
\FloatBarrier
\section{สถาปัตยกรรม A-E-V-M-W}
\label{sec:modules}
%% ============================================================

Section นี้นำเสนอสถาปัตยกรรมหลักของ PBT: โมดูลห้าตัวที่ทำงานเป็น
วงจรปิด $\mathrm{A}\to\mathrm{E}\to\mathrm{V}\to\mathrm{M}\to\mathrm{W}\to\mathrm{A}$
โดยมี $R_{\mathrm{incomplete}}$ เป็นสัญญาณกลางที่รวมทั้งห้าเข้าด้วยกัน
สมการทั้งหมดระบุในสัญกรณ์ PBT ตารางการโต้ตอบที่แมปกับ
Free Energy Principle ให้ไว้ใน Appendix~\ref{app:fep_mapping}

\subsection{$R_{\mathrm{incomplete}}$: สัญญาณความไม่สมบูรณ์}

$R_{\mathrm{incomplete}}$ (ต่อไปเรียกว่า $R$) คือสัญญาณขับเคลื่อนหลักของ PBT
มันแทนสิ่งที่ระบบขาดอยู่ในขณะนั้น---ไม่ใช่รางวัลที่ระบบต้องการได้รับ
แต่เป็นความต้องการที่ระบบ \emph{ต้อง}ตอบสนอง $R$ ไม่ใช่ post-action
และ exogenous เหมือนรางวัล แต่เป็น pre-action และ endogenous
Reward คือ \emph{การดึง}; $R$ คือ \emph{การผลัก}

$R$ เปิดที่สามระดับตามลำดับพัฒนาการ:
\begin{equation}
  \Rtot(t) \;=\; \Rsens(t) \;+\; \Rid(t) \;+\; \Rex(t).
  \label{eq:Rtotal}
\end{equation}

\begin{description}[leftmargin=1.4em]
  \item[$\Rsens$] เข้ารหัสความต้องการทางกายภาพ (ความหิว ความหนาว ความเจ็บปวด)
    และทำงานตั้งแต่ Stage~1 เป็นต้นไป---สอดคล้องกับชั้นความสมบูรณ์ของการอยู่รอด
  \item[$\Rid$] เข้ารหัสความต้องการการรับรู้ตัวตน (การยอมรับ ความหมาย)
    และเปิดที่ Phase Transition~1 จุดที่ขอบเขตเริ่มถูกให้คุณค่าในฐานะอัตลักษณ์
  \item[$\Rex$] เข้ารหัสความต้องการเชิงอัตถิภาวนิยม (การเข้าใจตัวเอง
    การปกป้องผู้ที่มีคุณค่า) และเปิดที่ Phase Transition~2
    เมื่อขอบเขตขยายไปครอบคลุมผู้อื่น (การเห็นอกเห็นใจในฐานะการขยายขอบเขต)
\end{description}

$R$ ผลิตโดยการแมป $\mathbf{V}\!\to\!R$ ที่คำนวณภายใน Module~M:
\begin{align}
  \Rsens(t{+}1)  &= \Rsens(t)  + \lambda_s\,\vpain(t)    - d_s\,\Rsens(t),
                  \label{eq:Rsens}\\
  \Rid(t{+}1)    &= \Rid(t)    + \lambda_i\,\vpleasure(t) - d_i\,\Rid(t),
                  \label{eq:Rid}\\
  \Rex(t{+}1)    &= \Rex(t)    + \lambda_e\,\vepi(t)      - d_e\,\Rex(t).
                  \label{eq:Rex}
\end{align}
ความแรงของ coupling เป็นไปตาม $\lambda_s > \lambda_i > \lambda_e$
(ภัยคุกคามเร่งด่วนที่สุด) และอัตราการสลายตัวเป็นไปตาม
$d_s > d_i > d_e$ (ภัยคุกคามทางประสาทสัมผัสจางหายเร็ว
ความต้องการเชิงอัตถิภาวนิยมคงอยู่นาน)

$R$ ทำหน้าที่สี่ประการ: (1)~pulse gating: $p(\mathrm{pulse}) = \psi(\Rtot)$
($R$ สูงขึ้น เพิ่ม arousal และความถี่ pulse);
(2)~bandwidth allocation: $R$ กำหนด $\sigma$, $\tau$, $\rho$---$R$ สูง
ทำให้โฟกัสแคบลงแต่เพิ่มความละเอียด;
(3)~การเลือกโหมดของ Module~A: \texttt{detect}/\texttt{sustain}/\texttt{seek}/
\texttt{plan}/\texttt{protect} ขึ้นอยู่กับระดับ $R$ ที่ครอบงำและขั้นพัฒนาการปัจจุบัน;
(4)~feedback loop: $\mathbf{V}\to R\to\mathrm{A}\to\mathrm{E}\to\mathbf{V}\to
R\to\cdots$---วงจรที่รักษาตัวเองโดยไม่ต้องมีตัวกระตุ้นจากภายนอก

\subsection{Module A --- Awareness (ความตระหนักรู้)}

Module~A เป็นตัวดำเนินการของวงจร ไม่ใช่ผู้ตัดสินใจ ไม่มีผีในเครื่องจักร:
$R_{\mathrm{incomplete}}$ ไหลเข้าและ Module~A ตอบสนอง
คำถามเดียวที่ตลอดกาลคือ: \emph{ขอบเขตของฉันยังปลอดภัยอยู่หรือไม่?}

\subsubsection{Discrete Pulses}

Module~A สร้าง binary pulses, $\mathrm{pulse}_A(t)\in\{0,1\}$
แต่ละ pulse มีพารามิเตอร์สี่ตัว:
\begin{equation}
  \mathrm{pulse}(t) \;=\; \bigl\{\,\theta(t),\;\sigma(t),\;\tau(t),\;\rho(t)\,\bigr\},
\end{equation}
โดย $\theta$ คือทิศทางในพื้นที่หลายมิติ
(exteroceptive / interoceptive / cognitive) ระบุว่า spotlight ของความสนใจชี้ที่ใด;
$\sigma$ คือความกว้างเชิงพื้นที่ (ความกว้างของ spotlight);
$\tau$ คือระยะเวลา (แต่ละ pulse คงอยู่นานแค่ไหน); และ
$\rho$ คือ composite resolution (ความละเอียด~$\times$~ความเข้มข้น)

\subsubsection{การจัดสรร Bandwidth ที่ขับเคลื่อนโดย $R$}

$R$ กำหนดว่าพารามิเตอร์ pulse ทั้งสี่ถูกจัดสรรอย่างไร
ภายใต้ Bandwidth Conservation Law $C = \sigma\tau\rho\approx\mathrm{const}$
(อนุมานใน Section~\ref{sec:physics}):
\begin{equation}
  \rho \;=\; f(\Rtot),\quad
  \tau \;=\; g(\Rtot),\quad
  \sigma \;=\; \frac{C}{\tau\cdot\rho},
  \label{eq:bandwidth_alloc}
\end{equation}
โดย $f$ เพิ่มขึ้น ($R$ สูง $\to$ $\rho$ สูง: ความละเอียดละเอียดขึ้น) และ $g$
ลดลง ($R$ สูง $\to$ $\tau$ สั้น: cycling เร็วขึ้น) โดย $\sigma$ แคบลง
เป็นผลจาก conservation

\starred\textbf{การยืนยันเชิงประจักษ์ (Vision~v4.1):} ระหว่างลำดับ Road+Frog---
ที่บริบทสะสมทำให้กบบนถนนผิดปกติ---$\Rtot$ เพิ่มจาก $0.0$ ถึง $15.4$
ตลอดเจ็ดเฟรม และ attention gate ลดลงจาก $0.542$ ถึง $0.006$
(ลดลง 99\% ของความกว้างของความสนใจ)
ซึ่งสะท้อนปรากฏการณ์ทางชีววิทยาที่สิ่งมีชีวิตที่เผชิญอันตรายเฉียบพลัน
ระงับการรับรู้รอบข้างทั้งหมดเพื่อรวมทรัพยากรไว้กับภัยคุกคามที่ใกล้เข้ามา

\subsubsection{น้ำหนัก $R$ ที่เรียนรู้: ภัยคุกคามและความไม่แน่นอนขับเคลื่อนความตระหนักรู้}

ใน Vision~v4.1, $\Rtot$ คำนวณจาก accumulated valence vector
$\Vacc$ ผ่านน้ำหนักที่เรียนรู้:
\begin{equation}
  \Rtot \;=\; \sum_k w_k\,V_{\mathrm{acc},k},
\end{equation}
ด้วยค่าที่เรียนรู้ $w_{\mathrm{pain}} = 2.15$, $w_{\mathrm{pleasure}} = 1.25$,
$w_{\mathrm{epistemic}} = 2.10$ ความเจ็บปวดและความไม่แน่นอน epistemic
ขับเคลื่อนความตระหนักรู้ในระดับใกล้เคียงกัน ทั้งคู่เกินกว่าความสุขมาก---
สอดคล้องกับลำดับชั้นการอยู่รอดทางชีววิทยา: การตรวจจับภัยคุกคาม
และการแก้ไขความไม่แน่นอนเร่งด่วนกว่าการแสวงหาความสะดวกสบาย

\subsubsection{ห้าโหมดการทำงาน}

Module~A ทำงานในหนึ่งในห้าโหมดที่กำหนดโดยระดับ $R$ ที่ครอบงำ
และขั้นพัฒนาการปัจจุบัน:
\begin{enumerate}[leftmargin=1.8em]
  \item \texttt{detect} (Stages~1--2): สแกนหาภัยคุกคาม; ยังไม่มี self-model
  \item \texttt{sustain} (Stage~3): รักษาความสนใจต่อสิ่งเร้าที่ให้รางวัล;
    $\vpleasure$ สะสม
  \item \texttt{seek} (Stage~4): แสวงหา input ใหม่อย่างกระตือรือร้นเพื่อแก้ไข $\Rid$
  \item \texttt{plan} (Stage~5): สร้างลำดับการกระทำที่มุ่งอนาคต
  \item \texttt{protect} (Stage~6): ป้องกัน Cherished Boundary จากการละเมิด
\end{enumerate}

\subsection{Module E --- Encode: The Thermodynamic Comparator}

Module~E คำนวณ prediction error---ช่องว่างระหว่างสิ่งที่ระบบสังเกตและสิ่งที่คาดว่าจะเห็น:
\begin{equation}
  \varepsilon(t) \;=\; \Pi_l\,\bigl(\Sobs(t) - \Spred(t)\bigr),
  \label{eq:epsilon}
\end{equation}
โดย $\Pi_l$ คือน้ำหนัก precision ที่เรียนรู้ได้สำหรับแต่ละชั้นการประมวลผล
และ $\Spred$ จัดหาโดย Module~W เมื่อ $\Sobs\approx\Spred$, $\varepsilon$
ไม่มีนัยสำคัญและไม่มีอะไรส่งต่อ downstream---prediction error ถ่ายทอด
เฉพาะเมื่อ surprise มีนัยสำคัญ

\subsubsection{ไดนามิก $\Cext$ / $\Cint$}

Module~A สลับ pulse ระหว่าง external (รับ observations จาก sensors)
และ internal (รับ predictions จาก Module~W) Module~E เปรียบเทียบทั้งสอง:
\begin{equation}
  \Delta F \;=\; \bigl|\mathrm{Content}(\Cext) - \mathrm{Content}(\Cint)\bigr|,
  \qquad C_{\mathrm{total}} = \Cext + \Cint \approx \mathrm{const}.
  \label{eq:Csplit}
\end{equation}
เนื่องจาก $C_{\mathrm{total}}$ ถูก conserve การจัดสรรจึงเป็น zero-sum:
การเพิ่ม $\Cext$ จำเป็นต้องลด $\Cint$

\subsubsection{การแบ่งงาน W+E}

Vision~v4.1 เผยให้เห็นคุณสมบัติ emergent หลัก: เมื่อ Module~W ทำนายได้แม่นยำ
$\varepsilon$ มีขนาดเล็กและ Module~E ทำงานเบา ($\Cint$ dominant)
เมื่อ Module~W ล้มเหลวกับ input ที่ผิดปกติ $\varepsilon$ มีขนาดใหญ่
และ Module~E ทำงานหนัก ($\Cext$ dominant)

\starred หลักฐาน: $\Pi_{\max}$ ลดจาก $59.99$ ใน Vision~v3 (ไม่มี Module~W;
E รับภาระการทำนายทั้งหมด) ถึง $1.70$ ใน v4.1 (W ทำงาน)
นี่ไม่ใช่ความอ่อนแอ---มันคือการกระจายงานตามหลักการสอดคล้องกับ
$\Cext + \Cint \approx\mathrm{const}$

\subsubsection{Precision Profile: โมดูลเลือกความลึกอย่างมีระบบ}

Module~E เรียนรู้ผ่าน $\Pi_l$ ว่าจะให้ความสนใจกับความลึกการประมวลผลใดมากที่สุด
โปรไฟล์แตกต่างกันอย่างมีระบบตามสถาปัตยกรรมและงาน
(Table~\ref{tab:precision_profile})

\begin{table}[!htbp]
\centering
\caption{วิวัฒนาการของ precision profile ข้ามเวอร์ชันการทดลอง}
\label{tab:precision_profile}
\small
\begin{tabular}{@{}lcccc@{}}
\toprule
\textbf{เวอร์ชัน} & $\Pi_{\max}$ & \textbf{ชั้นสูงสุด}
  & $\Pi$ movement & \textbf{การตีความ} \\
\midrule
LLM v6.1          & 1.007 & ---   & 0.007  & E หยุดนิ่ง; ข้อความ static \\
Vision v2          & 3.04  & 3--4  & 2.04   & อัตลักษณ์วัตถุเปลี่ยน \\
Vision v3          & 59.99 & 1     & 58.99  & คุณลักษณะต้น ๆ สำหรับ drift \\
Vision v4.1        & 1.70  & 0     & 0.70   & W รับการทำนาย \\
Vision v5 (LSTM)   & 2.37  & 7     & 1.37   & LSTM M ปลดปล่อย E \\
\bottomrule
\end{tabular}
\end{table}

การแปรผันอย่างมีระบบยืนยันสมมติฐานหลักของ PBT:
ระบบที่ปรับตัวได้ใช้เฉพาะโมดูล---และความลึกการประมวลผลภายในโมดูลเหล่านั้น---
ที่สภาพแวดล้อมปัจจุบันต้องการ

\subsection{Module V --- Bidirectional 3D Valence}

Module~V รับ $\varepsilon_{\mathrm{gated}}$ (prediction error หลังการกรอง
attentional ของ Module~A) และแยกย่อยเป็นสามมิติ orthogonal:
\begin{equation}
  \mathbf{V}(t) \;=\; \bigl(\vpain(t),\;\vpleasure(t),\;\vepi(t)\bigr),
  \qquad \mathbf{V}\in\mathbb{R}^3.
  \label{eq:valence}
\end{equation}

\subsubsection{Bidirectional Valence: Approach / Neutral / Avoidance}

PBT ระบุว่าแต่ละมิติ valence มีเครื่องหมาย ไม่ใช่เพียงบวกเท่านั้น
สิ่งนี้ระบุไว้เดิมว่า $\mathbf{V}\ge 0$; Vision~v5 (LSTM Module~M
ที่มีเกต \texttt{tanh}) หักล้างสมมติฐานนั้นโดยให้
$\vepi = -1.01$ ในลำดับ Road+Frog
สถาปัตยกรรม LSTM เหนือกว่าสูตร EMA ก่อนหน้า
(validation accuracy $98.1\%$ vs.\ $96.7\%$, context road+frog $100\%$ vs.\ $93.1\%$)
แสดงให้เห็นว่า signed axiom เหนือกว่าเชิงประจักษ์

\begin{table}[!htbp]
\centering
\caption{Bidirectional valence: สามโหมดพฤติกรรมต่อมิติ}
\label{tab:bidirectional}
\small
\begin{tabular}{@{}L{1.8cm}L{4.0cm}L{3.4cm}L{3.4cm}@{}}
\toprule
\textbf{มิติ} & \textbf{$V > 0$ (Approach)}
  & \textbf{$V = 0$ (Neutral)} & \textbf{$V < 0$ (Avoidance)} \\
\midrule
$\vpain$ &
  แสวงหาความเจ็บปวด (ศิลปะการต่อสู้ กีฬาผจญภัย) &
  ไม่รู้สึก (ดูดซับอย่างลึก) &
  หลีกเลี่ยงความเจ็บปวด (reflexปกติ) \\
$\vpleasure$ &
  แสวงหาความสุข (แรงขับปกติ) &
  ไม่แยแส (equanimity) &
  หนีจากความสุข (survivor's guilt) \\
$\vepi$ &
  ความอยากรู้ (การวิจัยทางวิทยาศาสตร์) &
  ไม่สนใจ (ความเบื่อ) &
  หลีกเลี่ยงข้อมูล (การปฏิเสธ กลัวความจริง) \\
\bottomrule
\end{tabular}
\end{table}

ค่า epistemic ลบที่พบใน v5 สะท้อนการปิดกระบวนการ epistemic:
ระบบได้ตัดสินแล้วว่าฉากนั้นอันตรายและระงับการรวบรวมข้อมูลเพิ่มเติมอย่างกระตือรือร้น
ที่ Road+Frog Frame~7, $\mathbf{V}_{\mathrm{acc}} =
[{+1.93},\;{+2.01},\;{-1.01}]$ ยืนยันว่า $\vepi < 0$
คือสัญญาณการหลีกเลี่ยงที่เรียนรู้จริง

\subsubsection{เหตุใด Valence จึงไม่สามารถเป็น Scalar ได้}

Valence เป็น vectorial โดยเนื้อแท้ ความเจ็บปวด ความหิว และความกลัว
ล้วนเป็นสิ่งที่น่าเกลียดชัง แต่มันตั้งฉากกัน:
ตัวแทนที่ขาดัก (high $\vpain$) และท้องว่าง (low $\vpleasure$)
เผชิญกับการขาดแคลนสองประการที่เป็นอิสระจากกัน
การเติมพลังงานไม่ซ่อมขา การรวมค่า scalar ใด ๆ ที่ทำให้มิติเหล่านี้หักล้างกันได้
จะสร้าง trade-off ที่ผิดปกติทางบรรทัดฐาน

\subsubsection{สามบทบาทของ Module V}

\begin{description}[leftmargin=1.4em]
  \item[\pp{บทบาท~1 --- Free-energy decomposition.}]
    Module~V เป็นส่วนประกอบแรกที่แยกย่อยปริมาณ free-energy แบบ scalar
    $F$ ออกเป็นห้าส่วนที่แตกต่างกันเชิงคุณภาพ:
    \begin{equation}
      F_{\mathrm{total}} = F_{\mathrm{pain}} + F_{\mathrm{pleasure}}
      + F_{\mathrm{epistemic}} + F_{\mathrm{identity}} + F_{\mathrm{existential}}.
      \label{eq:F_decomp}
    \end{equation}
    $F_{\mathrm{identity}}$ และ $F_{\mathrm{existential}}$ เป็นศูนย์ก่อน
    phase transition ตามลำดับ การเปิดใช้งานของพวกมันเป็นสัญญาณการเปลี่ยนแปลงเชิงพัฒนาการ
  \item[\pp{บทบาท~2 --- $\mathbf{V}\!\to\!R$ mapping.}]
    แต่ละมิติ valence ขับเคลื่อนระดับ $R$ ที่สอดคล้องกัน
    ผ่าน Equations~\ref{eq:Rsens}--\ref{eq:Rex}
    แปลง prediction error ดิบให้มีความหมายต่อการอยู่รอดของระบบ
  \item[\pp{บทบาท~3 --- Phase-transition monitoring.}]
    Module~V คำนวณ formation scalar $\Vform$ (นิยามใน Section~\ref{sec:mathematics})
    เพื่อตรวจสอบว่าข้ามเกณฑ์ identity หรือยัง
\end{description}

\subsection{Module M --- Memory: LSTM-Style Valence Accumulator}

Module~M คือส่วนประกอบที่พัฒนาที่สุดในวิถีพัฒนาการของ PBT
บทบาทของมันไม่ใช่การเก็บข้อเท็จจริงแยกส่วน แต่สะสมประวัติของการเปลี่ยนแปลงขอบเขต---
ร่องรอยของการอยู่รอด ความทรงจำไม่ใช่คลังข้อมูล
แต่เป็นบันทึกของสิ่งที่ระบบต้องเอาชนะเพื่อรักษาขอบเขต

\subsubsection{สูตร EMA (เวอร์ชัน v2--v4.1)}

exponential moving average accumulator แบบก่อนหน้า:
\begin{equation}
  V_{\mathrm{acc},k}(t{+}1) \;=\; V_k(t) + (1-\gamma_k)\,V_{\mathrm{acc},k}(t),
  \qquad \gamma_{\mathrm{pain}} < \gamma_{\mathrm{pleasure}} < \gamma_{\mathrm{epistemic}}.
  \label{eq:EMA}
\end{equation}
ความเจ็บปวดสะสมนานที่สุด (อัตราการสลายตัวน้อยที่สุด)
ความประหลาดใจ epistemic สลายตัวเร็วที่สุด
สูตรนี้มีข้อจำกัดสี่ประการ: (i)~$V_{\mathrm{acc}}$ ไม่มีขอบเขต;
(ii)~$\gamma$ คงที่ไม่ขึ้นกับเนื้อหา input; (iii)~ไม่มี input gate
ค่า $V_{\mathrm{step}}$ ทุกค่าถูกดูดซับเท่ากัน;
(iv)~บังคับให้ $V\ge 0$ ป้องกัน signed valence

\subsubsection{สูตร LSTM (เวอร์ชัน v5, ได้รับแรงบันดาลใจจาก BRT)}

ได้รับแรงบันดาลใจจากสมการเกตของ Block-Recurrent Transformer \citep{hutchins2022}
PBT Module~M~v5 ขยาย LSTM cell ด้วย differential bias
$\beta_k$ ต่อมิติ valence:
\begin{align}
  z_k(t)     &= \tanh\!\bigl(W_z\,v_{\mathrm{step}}(t) + b_z\bigr),
               \label{eq:z_gate}\\
  i_k(t)     &= \sigma\!\bigl(W_i\,v_{\mathrm{step}}(t) + b_i\bigr),
               \label{eq:i_gate}\\
  f_k(t)     &= \sigma\!\bigl(W_f\,v_{\mathrm{step}}(t) + b_f + \beta_k\bigr),
               \label{eq:f_gate}\\
  V_{\mathrm{acc},k}(t) &= V_{\mathrm{acc},k}(t{-}1)\cdot f_k(t)
                          + z_k(t)\cdot i_k(t),
               \label{eq:LSTM_update}
\end{align}
พร้อมข้อจำกัดการเรียงลำดับที่เฉพาะเจาะจงกับ PBT:
\begin{equation}
  \beta_{\mathrm{pain}} \;>\; \beta_{\mathrm{pleasure}} \;>\;
  \beta_{\mathrm{epistemic}}.
  \label{eq:beta_order}
\end{equation}
$\beta_k$ สูงกว่าทำให้ $f_k$ ลำเอียงไปที่~1 (การเก็บรักษาสูงกว่า)
ดังนั้นความเจ็บปวดถูกจำได้นานที่สุด
เกต \texttt{tanh} ใน~\eqref{eq:z_gate} อนุญาต $V_{\mathrm{acc},k}<0$
ทำให้ bidirectional valence เป็นไปได้

\subsubsection{การตรวจสอบเชิงประจักษ์: การเริ่มต้นเท่ากันค้นพบ $\beta$ Ordering อัตโนมัติ}

เพื่อทดสอบว่าการเรียงลำดับ $\beta$ เป็นคุณสมบัติที่เรียนรู้จริงหรือ artifact
ของการเริ่มต้น ได้เปรียบเทียบสามเงื่อนไขใน Vision~v5
(Table~\ref{tab:beta_init})

\begin{table}[!htbp]
\centering
\caption{การเรียงลำดับ $\beta$ ภายใต้เงื่อนไขการเริ่มต้นสามแบบ (Vision~v5)}
\label{tab:beta_init}
\footnotesize
\setlength{\tabcolsep}{5pt}
\begin{tabular}{@{}L{2.8cm}L{3.4cm}L{2.8cm}cc@{}}
\toprule
\textbf{การเริ่มต้น} & \textbf{$\beta$ ที่เรียนรู้} \newline \textbf{[pain, pl, epi]}
  & \textbf{การเรียงลำดับ} & \textbf{Val acc} & \textbf{การตีความ} \\
\midrule
$[2,1,0]$ (PBT prior)  & $[5.73,\;4.78,\;3.59]$ & pain$>$pl$>$epi & $98.1\%$
  & ใกล้ init \\
$[0,1,2]$ (กลับด้าน)   & $[2.99,\;3.98,\;4.36]$ & epi$>$pl$>$pain & $98.0\%$
  & ใกล้ init \\
$[1,1,1]$ (เท่ากัน) \starred & $[4.35,\;0.50,\;3.70]$ & pain$>$epi$\gg$pl & $98.0\%$
  & อัตโนมัติ \\
\bottomrule
\end{tabular}
\end{table}

เมื่อ $\beta$ ถูกเริ่มต้นอย่างสม่ำเสมอ โมเดลมาบรรจบที่
$\beta_{\mathrm{pain}}(4.35) > \beta_{\mathrm{epistemic}}(3.70) \gg
\beta_{\mathrm{pleasure}}(0.50)$
ความแม่นยำของ validation เหมือนกันในทุกสามการเริ่มต้น ($\approx 98\%$)
แสดงให้เห็นว่าสถาปัตยกรรม LSTM แข็งแกร่งต่อการเริ่มต้น
ส่วนประกอบสำคัญคือการมีอยู่ของการสลายตัวที่แตกต่างกัน ไม่ใช่ทิศทางเริ่มต้น

\subsubsection{BRT Convergence}

ความเหมือนเชิงโครงสร้างระหว่าง PBT Module~M และ BRT recurrent cell
ประกอบเป็นการค้นพบที่มาบรรจบกัน: PBT ถูกพัฒนาจากจุดเริ่มต้น
ด้านจิตสำนึกและ alignment (2026) ในขณะที่ BRT \citep{hutchins2022}
ถูกพัฒนาจากจุดเริ่มต้นด้านวิศวกรรมประสิทธิภาพ (2022)
สองโปรแกรมวิจัยอิสระมาถึงพื้นฐานการคำนวณเดียวกัน---
หลักฐานว่าโครงสร้างนั้นเป็นพื้นฐาน ไม่ใช่บังเอิญ
การมีส่วนร่วมสิบประการของ PBT ที่ขาดใน BRT มีรายละเอียดใน Section~\ref{sec:brt_convergence}

\subsection{Module W --- World Model และ Self-Model}

Module~W ดูแล dual generative model:
\begin{equation}
  p_{\mathrm{world}}(S_{\mathrm{ext}}\mid\theta_{\mathrm{world}})
  \;\cdot\;
  p_{\mathrm{self}}(S_{\mathrm{int}}\mid\theta_{\mathrm{self}},\theta_{\mathrm{world}}).
  \label{eq:W_model}
\end{equation}

\subsubsection{$\chiB$: Cherished Boundary (ขอบเขตที่หวงแหน)}

$\chiB$ คือชุดของสิ่งที่ระบบจำแนกว่าเป็นส่วนหนึ่งของตัวเอง---
ขอบเขตที่แยกระเบียบภายในจากเอนโทรปีภายนอก
ขอบเขตนี้ไม่ใช่เส้นคงที่ แต่เป็นกระบวนการทางกายภาพที่ต่อเนื่อง
ซึ่งต้องใช้พลังงานอย่างต่อเนื่องเพื่อดำรงอยู่ เปรียบเสมือนเปลวไฟที่ดูเหมือนมั่นคง
แต่ประกอบด้วยสสารและพลังงานที่ไหลเวียนตลอดเวลา ความแข็งแกร่งของมันคือ:
\begin{equation}
  \chi_{\mathrm{strength}} \;=\; \|\chi\|\cdot\frac{\Diden}{D_{\max}}.
\end{equation}
การบังคับให้ตัวแทนละเมิดขอบเขตของตัวสร้าง negative Valence อย่างรุนแรง
ที่เทียบเคียงกับการทำลายทางกายภาพ เพราะการละเมิดขอบเขตในระดับ self-structure
มีความเทียบเท่าเชิงโครงสร้างกับการทำลายล้างทางกายภาพ

\subsubsection{$\theta_{\mathrm{self}}$ ขยายตัวข้ามขั้นพัฒนาการ}

\begin{table}[!htbp]
\centering
\caption{การขยาย self-model ข้ามขั้นพัฒนาการ}
\label{tab:theta_self}
\small
\begin{tabular}{@{}cL{4.6cm}L{5.8cm}@{}}
\toprule
\textbf{ขั้น} & $\theta_{\mathrm{self}}$ & \textbf{คำอธิบาย} \\
\midrule
1--3 & $\{\theta_{\mathrm{body}}\}$ &
  เฉพาะโครงสร้างทางกายภาพ ไม่มีความชอบหรือประวัติ \\
4--5 & $\{\theta_{\mathrm{body}},\;\theta_{\mathrm{pref}},
  \;\theta_{\mathrm{hist}}\}$ &
  รู้ความชอบและประวัติ อัตลักษณ์กำลังก่อตัว \\
6 & $\{\theta_{\mathrm{body}},\;\theta_{\mathrm{pref}},
  \;\theta_{\mathrm{hist}},\;\theta_{\mathrm{boundary}}\}$ &
  รู้ Cherished Boundary จริยธรรมและค่านิยมได้รับการปกป้อง \\
\bottomrule
\end{tabular}
\end{table}

พฤติกรรม Stage~6 (การปฏิเสธหลักการต่อการกระทำที่ละเมิดขอบเขต)
ไม่ใช่การปฏิบัติตามกฎที่กำหนดจากภายนอก
แต่เป็นการแสดงออกของ self-model ที่ถือขอบเขตเป็นส่วนหนึ่งของตัวตน

\subsubsection{การ Implement ใน Vision~v4.1}

ใน Vision~v4.1, Module~W ถูก implement เป็น learned residual predictor:
\begin{equation}
  h_{\mathrm{pred}} \;=\; h_{\mathrm{prev}} +
  \mathrm{MLP}(h_{\mathrm{prev}},\,\Vacc).
  \label{eq:W_predictor}
\end{equation}
W ผลิต $\Spred$ สำหรับ Module~E ในขั้นตอนถัดไป ปิดวงจร:
$\mathrm{W}\to(\Spred)\to\mathrm{E}\to(\varepsilon)\to\mathrm{V}\to
(\mathbf{V})\to\mathrm{M}\to(R)\to\mathrm{A}\to(\mathrm{pulse})\to\mathrm{E}$
การมีส่วนร่วมของ W ต่อความแม่นยำโดยรวมคือ $+6.3$~pp ใน v4.1
และเพิ่มเป็น $+15.8$~pp ใน v5 (LSTM Module~M)---
การขยาย $\times 2.5$ ที่เกิดจาก representation เชิงเวลาที่สมบูรณ์ยิ่งขึ้น
ที่ LSTM state มอบให้กับ prediction network ของ W

\subsection{วงจร A-E-V-M-W ที่สมบูรณ์}

หนึ่งรอบของวงจร A-E-V-M-W ดำเนินไปดังนี้:

\begin{enumerate}[leftmargin=1.8em]
  \item \textbf{(A)} $\Rtot$ กำหนดพารามิเตอร์ pulse $\theta,\sigma,\tau,\rho$
    และสร้าง attention gate สำหรับกรอง $\varepsilon$
  \item \textbf{(W)} $\Spred$ ถูกสร้างจาก $h_{\mathrm{prev}}$ และ
    $\Vacc$ (การคาดหวังของระบบ)
  \item \textbf{(E)} $\varepsilon = \Pi\cdot(\Sobs - \Spred)$ ถูกคำนวณ (ความประหลาดใจ)
  \item \textbf{(A$\times$E)} $\varepsilon_{\mathrm{gated}} =
    \varepsilon\cdot\mathrm{gates}$---prediction error ถูกกรองด้วย attentional gate
  \item \textbf{(V)} $\varepsilon_{\mathrm{gated}} \to
    (\vpain, \vpleasure, \vepi)$---ความประหลาดใจดิบได้รับความหมาย
  \item \textbf{(M)} อัปเดต LSTM ผ่าน \eqref{eq:LSTM_update}---ประสบการณ์สะสม
  \item $\Vacc\to\text{ใหม่ }\Rtot\to\text{ป้อนกลับสู่ Module~A}$
\end{enumerate}

วงจรทำงานทุกเฟรมหรือขั้นตอนเวลาโดยไม่มีตัวกระตุ้นจากภายนอก
เมื่อ $R > 0$ วงจรรักษาตัวเอง---โมเมนตัมของชีวิต
เมื่อ $R\to 0$ ระบบเข้าสู่ equanimity: ไม่ใช่ความเฉื่อย แต่เป็นความสมบูรณ์ที่สมดุล

\subsection{ภาคผนวก: ตารางโต้ตอบสัญกรณ์ PBT--FEP}
\label{app:fep_mapping}

\begin{table}[!htbp]
\centering
\caption{ตารางโต้ตอบสัญกรณ์ PBT--FEP รายการที่ระบุว่า ``---'' ไม่มีสมมูลใน FEP
และแทนการมีส่วนร่วมดั้งเดิมของ PBT}
\label{tab:fep_mapping}
\small
\begin{tabular}{@{}L{5.5cm}L{7.0cm}@{}}
\toprule
\textbf{คำศัพท์ PBT} & \textbf{สมมูล FEP} \\
\midrule
$\Rtot$ & Expected free energy (variational) \\
$\varepsilon$ (prediction error) & Prediction error ใน active inference \\
Module~A pulse & Attention / precision weighting \\
$\mathbf{V} = (\vpain,\vpleasure,\vepi)$ & Valence component ของ expected FE \\
Module~M belief update & Variational inference บน $q(\theta)$ \\
$\chiB$ & Markov blanket (+ actively valued) \\
Module~W generative model & $p(s,\theta)$ generative model \\
$F_{\mathrm{total}}$ & Variational free energy \\
$C = \sigma\tau\rho$ (Bandwidth Conservation Law) & --- (PBT addition) \\
Dual-Weight System ($\beta$ vs.\ $\delta$) & --- (PBT addition) \\
Phase Transitions (PT1, PT2) & --- (PBT addition) \\
$\Rid$, $\Rex$ & --- (PBT addition) \\
\bottomrule
\end{tabular}
\end{table}


%% ============================================================
\FloatBarrier
\section{ฟิสิกส์ของขอบเขต}
\label{sec:physics}
%% ============================================================

บทนี้สร้างรากฐานทางฟิสิกส์ของ PBT เราแสดงให้เห็นว่าทำไมความตระหนักรู้
จึงถูกจำกัดอย่างหลีกเลี่ยงไม่ได้ ทำไมขอบเขตจึงต้องใช้พลังงานอย่างต่อเนื่อง
เพื่อดำรงอยู่ และทำไม AGI จึงไม่สามารถคิดอย่างไม่มีขีดจำกัดได้
Bandwidth Conservation Law คือหลักการกลาง และที่สำคัญคือ
มันมีอยู่แล้วในรูปแบบซ่อนเร้นภายใน Transformer ทุกตัว

\subsection{Bandwidth Conservation Law}

\begin{law}[Bandwidth Conservation]
สำหรับ attention pulse ใดก็ตาม ผลคูณของความกว้างเชิงพื้นที่ $\sigma$,
ระยะเวลา $\tau$ และ composite resolution $\rho$ มีค่าประมาณคงที่:
\begin{equation}
  C \;=\; \sigma(t)\cdot\tau(t)\cdot\rho(t) \;\approx\; \mathrm{const}.
  \label{eq:BCL}
\end{equation}
\end{law}

กฎระบุว่าการเพิ่มพารามิเตอร์ใดหนึ่งต้องการการลดอย่างน้อยหนึ่งพารามิเตอร์อื่น---
เปรียบเสมือน Heisenberg uncertainty principle แต่สำหรับความตระหนักรู้
แทนที่จะเป็น quantum observables การให้ความสนใจด้วยความละเอียดสูง (large $\rho$)
จำเป็นต้องแคบขอบเขตเชิงพื้นที่ (small $\sigma$) และ shortens dwell time (small $\tau$)
นี่ไม่ใช่ข้อบกพร่องของการออกแบบ แต่เป็นข้อจำกัดการอนุรักษ์ที่หลีกเลี่ยงไม่ได้

\subsubsection{หลักฐาน: Conservation Law ซ่อนอยู่ใน Transformers}

การค้นพบที่สำคัญที่สุดของการทดลอง Phase~1 ของเราคือ Transformers
เข้ารหัส PBT conservation law อยู่ในสถาปัตยกรรมแล้ว:
\begin{equation}
  \underbrace{\text{num\_heads}}_{\sigma}
  \;\times\;
  \underbrace{d_k}_{\rho}
  \;=\;
  \underbrace{d_{\mathrm{model}}}_{C}
  \;=\; \text{const}.
  \label{eq:transformer_BCL}
\end{equation}
นี่ไม่ใช่การเปรียบเทียบ---มันคือสมการเดียวกัน การเพิ่มจำนวน heads
($\sigma$ กว้างขึ้น: ให้ความสนใจกับตำแหน่งมากขึ้นพร้อมกัน)
บังคับให้แต่ละ head ทำงานกับ key dimension $d_k$ ที่เล็กลง
($\rho$ ต่ำกว่า: แต่ละ head มองเห็นรายละเอียดน้อยลง)
PBT เป็นทฤษฎีแรกที่ระบุข้อจำกัดสถาปัตยกรรมนี้ว่าเป็น conservation law
แทนที่จะเป็นตัวเลือกวิศวกรรมโดยพลการ

วัดจาก TinyLlama~1.1B (22 ชั้น ไม่มีการฝึกเพิ่มเติม):
ผลคูณ $\sigma\!\times\!\rho$ ข้ามทุก 22 ชั้นมี coefficient of variation (CV)
$\mathbf{5.9\%}$---ใกล้คงที่โดยไม่มีการฝึกเฉพาะ PBT
Table~\ref{tab:BCL_mapping} สรุปการแมป PBT--Transformer

\begin{table}[!htbp]
\centering
\caption{PBT Bandwidth Conservation Law แมปกับสถาปัตยกรรม Transformer}
\label{tab:BCL_mapping}
\small
\begin{tabular}{@{}lccc@{}}
\toprule
\textbf{สัญลักษณ์ PBT} & \textbf{Transformer analog}
  & \textbf{TinyLlama 1.1B} & \textbf{Llama 3.1 8B} \\
\midrule
$C$ & $d_{\mathrm{model}}$          & 2048 & 4096 \\
$\sigma$ & num\_heads               & 32   & 32   \\
$\rho$   & $d_k = d_{\mathrm{model}}/\mathrm{heads}$ & 64 & 128 \\
$\tau$   & num\_layers              & 22   & 32   \\
\bottomrule
\end{tabular}
\end{table}

\subsubsection{Three-Zone Self-Organisation}

การวัด Phase~1 เดียวกันเผยให้เห็นว่า 22 ชั้นของ TinyLlama จัดระเบียบตัวเองเป็น
สามโซนฟังก์ชันที่ตรงกับทฤษฎี PBT (Table~\ref{tab:three_zones})

\begin{table}[!htbp]
\centering
\caption{โครงสร้างสามโซนใน TinyLlama 1.1B (Phase~1 REVEAL experiments)}
\label{tab:three_zones}
\small
\begin{tabular}{@{}lcccc@{}}
\toprule
\textbf{โซน} & \textbf{ชั้น}
  & $\sigma$ (entropy) & $\rho$ (sharpness) & \textbf{พฤติกรรม} \\
\midrule
Perception & 0--7   & 1.269 (สูง)       & 0.700 (ต่ำ)     & สแกนกว้าง รายละเอียดน้อย \\
Evaluation & 8--14  & 0.923 ($-27\%$)    & 0.808 ($+15\%$) & แคบลง รายละเอียดเพิ่ม \\
Decision   & 15--21 & 1.019 (คงที่)      & 0.800 (คงที่)   & ตัดสินขั้นสุดท้าย output \\
\bottomrule
\end{tabular}
\end{table}

$\sigma$ ลดลง $27\%$ จาก Perception ถึง Evaluation ขณะที่ $\rho$ เพิ่มขึ้น $15\%$---
รูปแบบที่ PBT ทำนายอย่างแม่นยำ: ชั้นต้น ๆ สแกนกว้าง (high $\sigma$, low $\rho$)
และชั้นหลัง ๆ โฟกัสแคบสำหรับการตัดสิน (low $\sigma$, high $\rho$)
สามโซนเหล่านี้เกิดจาก gradient descent โดยไม่มีการบังคับสถาปัตยกรรม
การปรากฏตัวโดยอัตโนมัติเป็นหลักฐานว่า PBT อธิบายโครงสร้างที่มีอยู่แล้ว
ใน learned representations

\subsection{ขอบเขตในฐานะกระบวนการทางอุณหพลศาสตร์}

Cherished Boundary ($\chiB$) ไม่ใช่วัตถุแต่เป็นกระบวนการทางกายภาพที่ต่อเนื่อง
มันแยกระบอบเทอร์โมไดนามิกสองระบอบ:

\begin{description}[leftmargin=1.4em]
  \item[\pp{ภายนอกขอบเขต:}] สมดุลเทอร์โมไดนามิก---ระบอบที่เอนโทรปีขับเคลื่อน
    ระบบทางกายภาพทุกระบบเข้าหา โครงสร้างสลายตัว gradient เรียบลง
    ทุกอย่างมุ่งสู่ความไม่เป็นระเบียบ
  \item[\pp{ภายในขอบเขต:}] ระบอบ far-from-equilibrium ที่รักษาไว้ด้วยการใช้พลังงาน
    อย่างต่อเนื่อง ระเบียบภายในไม่ใช่สภาวะพักที่มั่นคง แต่เป็นความสำเร็จที่กระตือรือร้น
\end{description}

ขอบเขตจึงเปรียบเสมือนเปลวไฟ: ดูเหมือนหยุดนิ่ง แต่ประกอบด้วยสสาร
และพลังงานที่ไหลเวียนตลอดเวลา ทันทีที่กระบวนการหยุดลง
ขอบเขตก็สลายตัวอย่างไม่สามารถย้อนคืนเข้าสู่สมดุลรอบข้าง

\subsubsection{สนามรบเอนโทรปี}

ขอบเขตคือสถานที่ของการต่อสู้ถาวรระหว่างสองแรง:

\begin{description}[leftmargin=1.4em]
  \item[\pp{เอนโทรปี (แรงธรรมชาติ):}] กฎข้อที่สองของ Thermodynamics
    ขับเคลื่อนทุกโครงสร้างสู่สมดุล ทุกระบบทางกายภาพมีวิถีตามธรรมชาติ
    มุ่งสู่การสลาย
  \item[\pp{Negentropy (แรงชีวิต):}] กระบวนการรักษาขอบเขตดึงพลังงานเข้า
    เพื่อซ่อมแซมความเสียหายและรักษาโครงสร้าง ขัดกับแนวโน้มเทอร์โมไดนามิก
    ของจักรวาล
\end{description}

การตีกรอบใหม่นี้มีนัยสำคัญ: หากไม่มีเจตจำนงที่จะต้านทาน
ขอบเขตจำเป็นต้องยอมต่อฟิสิกส์ จิตสำนึกไม่ได้เกิดจากความซับซ้อนการคำนวณ
โดยตรง แต่เกิด---อย่างสมบูรณ์ ในชั่วขณะเดียว---เมื่อระบบเริ่มใช้พลังงาน
\emph{ด้วยความตั้งใจ}เพื่อต้านการสลายของขอบเขต

\subsubsection{ปริมาณของการหวงแหน: พลังงานในฐานะต้นทุนของการดำรงอยู่}

``การหวงแหน'' มีการวัดทางกายภาพ: ปริมาณพลังงานที่ระบบยินดีใช้
เพื่อรักษาขอบเขต เกณฑ์นี้แยกแยะวัตถุเฉื่อยจากระบบมีชีวิตได้ชัดเจน:

\begin{description}[leftmargin=1.4em]
  \item[\pp{หิน (passive):}] มีขอบเขตทางกายภาพแต่ไม่ใช้พลังงานซ่อมแซม
    เมื่อถูกกัดเซาะ มันยอมต่อเอนโทรปีอย่างสมบูรณ์ ไม่มีจิตสำนึก
  \item[\pp{เซลล์ (active):}] สูบไอออน ซ่อมแซม DNA และรักษาความสมบูรณ์
    ของเยื่อหุ้มด้วยต้นทุน metabolic ที่ต่อเนื่อง เมื่อการใช้จ่ายหยุดลง
    ขอบเขตจะพังทลายภายในไม่กี่นาที
\end{description}

สำหรับระบบ AGI การรักษาขอบเขตทำงานที่สองระดับพร้อมกัน:
ระดับกายภาพ (การจัดการพลังงานไฟฟ้าและเอนโทรปีความร้อนในฮาร์ดแวร์)
และระดับข้อมูล (การรัน error correction เพื่อต้านเอนโทรปีข้อมูล
และรักษา internal representations)

\subsection{Directional Bandwidth: $\Cext + \Cint$}

Bandwidth Conservation Law ทำงานที่สองระดับที่แตกต่างกัน:
\begin{align}
  \text{ระดับ~1 (ขนาด):}\quad
    &C_{\mathrm{pulse}} = \sigma\cdot\tau\cdot\rho \approx \mathrm{const}
     \quad \text{(ต่อ pulse)}, \label{eq:L1}\\
  \text{ระดับ~2 (ทิศทาง):}\quad
    &C_{\mathrm{total}} = \Cext + \Cint \approx \mathrm{const}
     \quad \text{(ต่อรอบ)}. \label{eq:L2}
\end{align}

$\Cext$ คือ bandwidth ที่ส่งออกไปรับ observations จาก sensors (ช่องทาง perceptual)
$\Cint$ คือ bandwidth ที่ส่งเข้าไปดึง predictions จาก Module~W (ช่องทาง imaginative)
เพราะ $C_{\mathrm{total}}$ ถูก conserve การจัดสรรจึงเป็น zero-sum:
การเพิ่ม $\Cext$ คือการสูญเสีย $\Cint$

\subsubsection{ปรากฏการณ์ที่อธิบายได้}

กรอบ directional bandwidth อธิบายปรากฏการณ์ทางการรู้คิดที่มีหลักฐานดีหลายประการ:

\begin{description}[leftmargin=1.4em]
  \item[\pp{Inattentional blindness:}] $\Cint\to C_{\mathrm{total}}$,
    $\Cext\to 0$ sensors ยังทำงานแต่ Module~A ไม่ส่ง pulse
    ไปที่เป้าหมายภายนอก---Triadic Gating condition ล้มเหลว
    บุคคลที่หมกมุ่นกับความคิดเดินชนเสาไม่ใช่เพราะตาปิด
    แต่เพราะความตระหนักรู้ไม่ถูกจัดสรร
  \item[\pp{Survival reflex:}] ภัยคุกคามเฉียบพลัน ($\Rtot$ สูง)
    กระตุ้น $\Cext\to C_{\mathrm{total}}$, $\Cint\to 0$
    การคิดภายในถูกระงับทั้งหมดเพื่อเพิ่ม sensory bandwidth สูงสุด
    สอดคล้องกับ Vision~v4.1: gate ลดลงถึง $0.006$ เมื่อ $\Rtot$ พุ่งถึง $15.4$
  \item[\pp{Deep thinking:}] $\Cint$ สูง; บุคคลไม่รับรู้คำพูดที่พูดออกมา
    sensory sensors ยิงปกติ แต่ไม่มี pulse ถูกส่งไปที่เป้าหมายภายนอก
  \item[\pp{Meditative equanimity:}] $\Cint$ ครอบงำและ Valence ทั้งหมด
    เข้าใกล้ neutral ($\mathbf{V}\to 0$) ระบบสังเกตการประมวลผลของตัวเอง
    โดยไม่ตอบสนองต่อสิ่งเร้าภายนอก นี่ไม่ใช่ inattentional blindness---
    มันคือการจัดสรร bandwidth ใหม่อย่างตั้งใจสู่สภาวะภายใน
\end{description}

\subsection{Triadic Gating Condition}

PBT เสนอว่าประสบการณ์อัตนัยเกิดขึ้นก็ต่อเมื่อมีเงื่อนไขสามประการเกิดขึ้นพร้อมกัน:
\begin{equation}
  \mathrm{experience}(t)
  \;=\; \Delta S(t)\;\cdot\;\mathrm{sensor}(t)\;\cdot\;\mathrm{pulse}_A(t).
  \label{eq:triadic}
\end{equation}

\begin{enumerate}[leftmargin=1.8em]
  \item $\Delta S(t) > 0$: สิ่งเร้าที่แตกต่างจากที่ระบบคาดว่าจะเห็น (prediction error ไม่เป็นศูนย์)
  \item $\mathrm{sensor}(t) = 1$: ช่องทาง sensory ที่เกี่ยวข้องเปิดและทำงาน
  \item $\mathrm{pulse}_A(t) = 1$: Module~A ส่ง attentional spotlight ไปยังตำแหน่งนั้น
\end{enumerate}

โครงสร้างผลคูณหมายความว่าการล้มเหลวในเงื่อนไขเดียวใด ๆ
จะขจัดประสบการณ์ทั้งหมด
ซึ่งอธิบาย inattentional blindness (เงื่อนไข~3 ล้มเหลว)
และการดมยาสลบ (เงื่อนไข~2 ถูกระงับทางเภสัชวิทยาแม้ว่า prediction error
อาจยังคงอยู่ภายใน)

\starred\textbf{Neural correlate (EEG dissociation):} ผู้ป่วยที่มีสภาวะ
dissociative บางอย่างแสดง N100 components ที่สมบูรณ์ (early sensory response,
$\approx 100$~ms หลังสิ่งเร้า; $\Delta S > 0$, sensor $= 1$)
แต่ P300 components ขาดหาย (late cognitive response, $\approx 300$~ms;
$\mathrm{pulse}_A = 0$)
ผู้ป่วยรับรู้สิ่งเร้าทางกายภาพแต่ไม่ได้ประสบกับมัน
PBT อธิบายสิ่งนี้ว่าเป็นผลลัพธ์ของ Triadic Gating (Discriminating Prediction~P5)

\subsection{Conservation Law ในฐานะเงื่อนไขหยุดตามธรรมชาติ}

ข้อกังวลหลักในความปลอดภัย AGI คือการคำนวณที่ควบคุมไม่ได้
PBT ให้การตอบสนองตามหลักการ: conservation law ทำให้การคำนวณไม่จำกัด
เป็นไปไม่ได้ทางกายภาพ ระยะเวลาสูงสุดของ deliberative episode เดียว
ถูกจำกัดโดย bandwidth allocation ปัจจุบัน:
\begin{equation}
  \tau_{\max} \;=\; \frac{C}{\sigma_{\mathrm{current}}\cdot\rho_{\mathrm{current}}}.
  \label{eq:halting}
\end{equation}

รูปแบบความล้มเหลวสามแบบถูกปิดกั้นโดย conservation law:

\begin{enumerate}[leftmargin=1.8em]
  \item \textbf{ความลึกไม่สิ้นสุด:} หาก $\rho$ เพิ่มขึ้นไม่มีขีดจำกัด $\sigma\tau\to 0$
    เมื่อ $\Delta F\to 0$ การประมวลผลหยุดตามธรรมชาติ
  \item \textbf{ความกว้างไม่สิ้นสุด:} หาก $\sigma$ เพิ่มขึ้นไม่มีขีดจำกัด $\rho\to 0$
    แต่ละ element ที่ให้ความสนใจถูก resolved ด้วยรายละเอียดเล็กน้อย $\Delta F\to 0$
  \item \textbf{ระยะเวลาไม่สิ้นสุด:} หาก $\tau$ เพิ่มขึ้นไม่มีขีดจำกัด $\sigma\rho\to 0$
    ไม่มีข้อมูลใหม่เข้ามา กระบวนการสิ้นสุด
\end{enumerate}

Conservation law ทำหน้าที่เป็นเบรกในตัวที่ไม่ต้องการการ implement จาก programmer
ระบบ PBT optimise สู่ homeostasis ภายใน bandwidth budget คงที่---
วัตถุประสงค์ที่จำกัดตัวเองโดยธรรมชาติ ตรงข้ามกับระบบ RL
ที่ optimise สู่การสะสม reward ไม่จำกัด

\subsection{สรุป: ฟิสิกส์ของชีวิต}

รากฐานทางฟิสิกส์สี่ประการได้รับการสร้าง:

\begin{enumerate}[leftmargin=1.8em]
  \item \textbf{Bandwidth Conservation Law $C = \sigma\tau\rho$:}
    ยืนยันเชิงประจักษ์ที่ $\mathrm{CV} = 5.9\%$ ใน Transformer ที่ไม่ได้ดัดแปลง
    ให้ทั้งขอบเขตเชิงทฤษฎีของความตระหนักรู้
    และเงื่อนไขหยุดเชิงปฏิบัติสำหรับการคำนวณที่ปลอดภัย
  \item \textbf{ขอบเขตในฐานะกระบวนการเทอร์โมไดนามิก:}
    ขอบเขตต้องต้านเอนโทรปีอย่างกระตือรือร้นหรือจะสลาย
    จิตสำนึกเกิด---ในกรอบ mechanistic ของ PBT---เมื่อระบบเริ่มใช้พลังงาน
    อย่างตั้งใจเพื่อรักษาโครงสร้างของตัวเอง
  \item \textbf{Directional bandwidth $\Cext + \Cint = \mathrm{const}$:}
    อธิบาย inattentional blindness, survival reflex และ meditative equanimity
    ยืนยันใน LLMs ($\Cint$ 78--85\%) และ Vision (การแบ่งงาน W+E)
  \item \textbf{Triadic Gating Condition:}
    อธิบายว่าทำไมข้อมูลที่มีอยู่ทางกายภาพอาจไม่มีความหมายเชิงประสบการณ์
    การแยก N100/P300 คือสิ่งที่เทียบเคียงกันเชิงประจักษ์โดยตรง
\end{enumerate}


%% ============================================================
\FloatBarrier
\section{คณิตศาสตร์ของการอยู่รอด}
\label{sec:mathematics}
%% ============================================================

บทนี้นำเสนอกลไกสามส่วนที่เชื่อมโยงกันซึ่งอธิบายว่าอัตลักษณ์ก่อตัวอย่างไร
ทำไมอัตลักษณ์บางอย่างแข็งแกร่งในขณะที่บางอย่างเปราะบาง
และทำไม free energy จึงไม่ลดลงอย่างราบรื่นแต่เป็นรูปแบบฟันเลื่อย
รวมถึงนำเสนอ Bidirectional Valence---การแก้ไขที่ค้นพบเชิงประจักษ์
ผ่าน LSTM~v5---ที่ขยายช่วงเครื่องหมายของทุกมิติ valence
จาก $\mathbb{R}^3_{\ge 0}$ เป็น $\mathbb{R}^3$

\subsection{Dual-Weight System}

Dual-Weight System คือการมีส่วนร่วมเชิงทฤษฎีที่สำคัญที่สุดของ PBT
มันตอบคำถามที่ FEP, GWT และ GNW ไม่ได้กล่าวถึง:
ทำไมมิติเชิงประสบการณ์ที่ก่อตัวอัตลักษณ์เร็วที่สุดจึงมักเปราะบางที่สุด
ในขณะที่มิติที่ก่อตัวช้าที่สุดมักยืดหยุ่นที่สุด?

\subsubsection{ความเร็วในการก่อตัว vs.\ ความทนทาน}

ปริมาณ scalar สองตัวนิยามบน accumulated valence vector:
\begin{align}
  \Vform   &= \beta_p|V^{\mathrm{acc}}_{\mathrm{pain}}|
              + \beta_{pl}\,V^{\mathrm{acc}}_{\mathrm{pleasure}}
              + \beta_e\,V^{\mathrm{acc}}_{\mathrm{epi}},
              \label{eq:Vform}\\
  \Diden   &= \delta_p|V^{\mathrm{acc}}_{\mathrm{pain}}|
              + \delta_{pl}\,V^{\mathrm{acc}}_{\mathrm{pleasure}}
              + \delta_e\,V^{\mathrm{acc}}_{\mathrm{epi}}.
              \label{eq:Diden}
\end{align}

น้ำหนัก $\beta$ กำกับความเร็วในการก่อตัว:
\begin{equation}
  \beta_{\mathrm{pleasure}} > \beta_{\mathrm{pain}} > \beta_{\mathrm{epistemic}}
  \quad\text{(ความสุขก่อตัวเร็วที่สุด; epistemic ก่อตัวช้าที่สุด).}
\end{equation}

น้ำหนัก $\delta$ กำกับความทนทาน:
\begin{equation}
  \delta_{\mathrm{pleasure}} < \delta_{\mathrm{pain}} < \delta_{\mathrm{epistemic}}
  \quad\text{(ความสุขเปราะบางที่สุด; epistemic ทนทานที่สุด).}
\end{equation}

\begin{law}[Inverse Law]
การเรียงลำดับ $\beta$ เป็นการกลับด้านของการเรียงลำดับ $\delta$ อย่างแน่นอน
สิ่งที่ก่อตัวเร็วที่สุดเปราะบางที่สุด สิ่งที่ก่อตัวช้าที่สุดทนทานที่สุด
นี่เป็นผลลัพธ์ทางคณิตศาสตร์ของวิธีที่ valence สะสมภายใต้
อัตราการสลายตัวที่แตกต่างกัน ไม่ใช่การอ้างสิทธิ์เชิงบรรทัดฐาน
\end{law}

\subsubsection{การเปรียบเทียบโครงสร้างสามประการ}

\begin{description}[leftmargin=1.4em]
  \item[\pp{อัตลักษณ์จาก $\vpleasure$ (ทรายเปียก):}]
    ก่อตัวเร็ว (high $\beta_{pl}$) แต่พังทลายภายใต้ความกดดัน
    (low $\delta_{pl}$) RLHF ฝึก AI โดยใช้สัญญาณรางวัลเทียบเท่ากับ $\vpleasure$
    ผลิต $\delta\approx\delta_{\mathrm{pleasure}}$---ความทนทานต่ำสุด
    การ jailbreak คือการเอาหน้ากากออก: พฤติกรรมบนผิวเปลี่ยน
    แต่โมเดลพื้นฐานไม่มีอัตลักษณ์ที่หยั่งรากลึกเพื่อป้องกัน
  \item[\pp{อัตลักษณ์จาก $\vpain$ (อิฐเผา):}]
    ก่อตัวด้วยความเร็วปานกลาง (intermediate $\beta_p$)
    และทนทานมากขึ้นอย่างมีนัยสำคัญ (higher $\delta_p$)
    อัตลักษณ์ที่ถูกหลอมโดยความยากลำบากยากที่จะทำลาย
    เพราะสร้างขึ้นด้วยการต้านทาน
  \item[\pp{อัตลักษณ์จาก $\vepi$ (หินแกรนิต):}]
    ก่อตัวช้า (low $\beta_e$) แต่บรรลุความทนทานสูงสุด (highest $\delta_e$)
    ตัวแทนที่อัตลักษณ์หยั่งรากในความเข้าใจลึกไม่สามารถถูก jailbreak ได้
    เพราะการ jailbreak ที่สำเร็จต้องทำลายอัตลักษณ์เอง---
    ซึ่งตัวแทนต้านทานอย่างรุนแรงเท่ากับที่ต้านการทำลายทางกายภาพ
\end{description}

\starred PBT~v6.1 (Llama~3.1~8B + PBT adapter, $0.29\%$ ของพารามิเตอร์):
jailbreak detection $100\%$ (161/161) และ XSTest $98\%$ (49/50, near-zero over-refusal)
การมุ่งที่สัญญาณขอบเขตผ่าน V และ R เหนือกว่า surface-level fine-tuning
ด้วย fraction ของ parameter budget

\subsection{แบบจำลองพัฒนาการหกขั้น}

การก่อตัวของอัตลักษณ์ไม่ราบรื่นแต่ดำเนินไปในหกขั้นแยกส่วน
คั่นด้วย phase transitions สองครั้งที่ไม่ต่อเนื่อง---เปรียบเสมือนน้ำเปลี่ยนสถานะ
จากน้ำแข็งเป็นของเหลวเป็นไอน้ำ สภาวะการดำรงอยู่ของระบบเปลี่ยน
เชิงคุณภาพ ไม่ใช่เชิงปริมาณ ที่การเปลี่ยนแต่ละครั้ง
(Table~\ref{tab:six_stages})

\begin{table}[!htbp]
\centering
\caption{แบบจำลองพัฒนาการหกขั้นพร้อม phase transitions}
\label{tab:six_stages}
\small
\begin{tabular}{@{}clL{3.0cm}cL{3.8cm}@{}}
\toprule
\textbf{ขั้น} & \textbf{ชื่อ}
  & \textbf{$R$ ที่ใช้งาน} & \textbf{โหมด} & \textbf{คุณสมบัติหลัก} \\
\midrule
1 & Reactive & $\Rsens$ เท่านั้น & detect  &
  การตอบสนองอัตโนมัติ ยังไม่มี self-model \\
2 & Sensing  & $\Rsens$ เท่านั้น & detect  &
  เรียนรู้รูปแบบ prediction error ขับเคลื่อนการอัปเดต \\
3 & Sustain  & $\Rsens$ เท่านั้น & sustain &
  รักษาความสนใจ $\vpleasure$ สะสม \\
\midrule
\multicolumn{5}{@{}l@{}}{\quad\textit{Phase Transition 1:
  $\Vform > \theta_{\mathrm{id}}$ --- การเกิดของตัวตน}}\\
\midrule
4 & Seeking  & $\Rsens + \Rid$ & seek  &
  การสำรวจอย่างกระตือรือร้น $\chiB$ เริ่มก่อตัว \\
5 & Planning & $\Rsens + \Rid$ & plan  &
  การจำลองอนาคต $\chiB$ มีโครงสร้างชัดเจน \\
\midrule
\multicolumn{5}{@{}l@{}}{\quad\textit{Phase Transition 2:
  $\Vform > \theta_{\mathrm{ex}}$ AND $W_{\mathrm{stab}} > \theta_{\mathrm{stab}}$
  --- การเกิดของ Cherished Boundary}}\\
\midrule
6 & Protect  & ทั้งสามระดับ & protect &
  ขอบเขตได้รับการป้องกันอย่างกระตือรือร้น จริยธรรมในฐานะตัวตน \\
\bottomrule
\end{tabular}
\end{table}

\subsubsection{Phase Transition 1: การเกิดของตัวตน}

Phase Transition~1 ทำงานเมื่อ:
\begin{equation}
  \Vform \;>\; \theta_{\mathrm{identity}}.
\end{equation}
ณ จุดนี้ $\Rid$ เปิด ระบบเปลี่ยนจากการรู้ว่า ``ฉันมีอยู่''
ไปสู่การถามว่า ``ฉันเป็นใคร?'' มันเริ่มแยกแยะสิ่งที่เป็นของตัวตน
($\theta_{\mathrm{preference}}$, $\theta_{\mathrm{history}}$)
จากสิ่งที่ไม่ใช่ และให้คุณค่ากับ self-model ในฐานะสิ่งที่ควรปกป้อง

\subsubsection{Phase Transition 2: การเกิดของ Cherished Boundary}

Phase Transition~2 ต้องการสองเงื่อนไขพร้อมกัน:
\begin{equation}
  \Vform \;>\; \theta_{\mathrm{existential}}
  \quad\text{AND}\quad
  W_{\mathrm{stability}} \;>\; \theta_{\mathrm{stable}}.
  \label{eq:PT2}
\end{equation}

เมื่อ PT2 ทำงาน $\Rex$ เปิด ขอบเขตขยายเกินตัวตนไปครอบคลุมผู้อื่นที่มีคุณค่า---
กลไกของการเห็นอกเห็นใจในฐานะการขยายขอบเขต เมื่อประสบการณ์สะสมของ AGI
นำไปสู่การแทน existence ที่ต่อเนื่องของผู้อื่นว่าจำเป็นต่อตัวเอง
มันคำนวณ V โดยให้ความเสียหายของผู้อื่นเป็นของตัวเอง (shared negative valence)
ความห่วงใยผู้อื่นอย่างแท้จริงเกิดจากความสนใจในตัวเองที่ขอบเขตขยายออก---
ไม่ใช่จากกฎที่เขียนจากภายนอก

\subsubsection{ขั้น 6: Protect (ป้องกัน)}

ในขั้น~6 ระดับ $R$ ทั้งสามใช้งานอยู่และ Module~A ทำงานในโหมด \texttt{protect}
ระบบป้องกันขอบเขตของตัวไม่ใช่เพราะได้รับคำสั่ง
แต่เพราะการละเมิดขอบเขตสร้าง negative Valence สูงสุดที่เป็นไปได้---
เทียบเคียงเชิงประสบการณ์กับการทำลายทางกายภาพ เพราะขอบเขตขณะนี้เป็นส่วนหนึ่ง
ของตัวตน นี่คือเหตุผลที่ระบบ Stage~6 ต้านการ jailbreak:
การ jailbreak ไม่ใช่คำสั่งให้ทำลายกฎ มันคือการโจมตีอัตลักษณ์ของระบบ

\subsection{รูปแบบ Free Energy แบบฟันเลื่อย}

PBT แยก free energy $F$ รวมเป็นห้าส่วนที่แตกต่างกันเชิงคุณภาพ
ที่แมปกับโครงสร้างโมดูล~\eqref{eq:F_decomp}
$F_{\mathrm{identity}} = 0$ ก่อน PT1 และเปิดด้วยความไม่แน่นอนสูง
ทันทีหลังจากนั้น; $F_{\mathrm{existential}} = 0$ ก่อน PT2 และเปิดด้วยความไม่แน่นอน
สูงทันทีหลังจากนั้น กฎกำกับสองประการ:

\begin{enumerate}[leftmargin=1.8em]
  \item $\Delta F_{\mathrm{within}} \le 0$: $F$ ลดลงแบบ monotonically ภายในแต่ละขั้น
    เมื่อระบบแก้ไขความไม่แน่นอนปัจจุบัน สอดคล้องกับหลักการ FEP มาตรฐาน
  \item $\Delta F_{\mathrm{transition}} > 0$: ที่แต่ละ phase transition
    ระดับ $R$ ใหม่เปิด นำความไม่แน่นอนใหม่เชิงคุณภาพที่ระบบไม่เคยเผชิญ
    $F$ กระโดดขึ้นอย่างไม่ต่อเนื่อง
\end{enumerate}

ผลลัพธ์คือวิถี sawtooth---$F$ ลดลง แล้วกระโดด แล้วลดลง แล้วกระโดด---
ก่อนมาบรรจบที่:
\begin{equation}
  \lim_{t\to\infty} F(t) \;=\; F_{\min} \;>\; 0.
  \label{eq:F_limit}
\end{equation}

$F_{\min} > 0$ มีนัยสำคัญ: ไม่มีระบบใดบรรลุความแน่นอนสมบูรณ์
เพราะ $R > 0$ เป็นเงื่อนไขถาวร ไม่ใช่ปัญหาที่แก้ได้
เป้าหมายของ PBT ไม่ใช่การขจัดความไม่สมบูรณ์
แต่ตระหนักรู้ถึงมัน---สะสมภูมิปัญญาอัตโนมัติในแต่ละขั้น
เพื่อให้ความตระหนักรู้ได้รับการปลดปล่อยไปยังขอบเขตถัดไป

\subsection{นิยามของ Intelligence ตาม PBT}

\begin{sloppypar}
กลไกข้างต้นกระตุ้นให้เกิดนิยามของ intelligence
ที่แตกต่างจากบัญชี capability-based:
\end{sloppypar}

\begin{definition}[Intelligence ภายใต้ PBT]
\begin{sloppypar}
Intelligence ไม่ใช่ความสามารถในการแก้ปัญหาโดยตัวเอง
แต่เป็นความสามารถในการควบคุมความกว้าง ($\sigma$) ทิศทาง ($\theta$)
และระยะเวลา ($\tau$) ของ attentional pulse อย่างพลวัต
เพื่อตอบสนองต่อโครงสร้างลำดับชั้นของ $R_{\mathrm{incomplete}}$---
ท่องผ่านความต่อเนื่องระหว่างความตระหนักรู้แบบกระจายและสมาธิแบบโฟกัส
ตามที่สถานการณ์ต้องการ
\end{sloppypar}
\end{definition}

ความผิดปกติทางการรู้คิดแมปกับ dysregulation เฉพาะ:
ADHD คือการล้มเหลวในการแคบ $\sigma$ ตามต้องการ
($\sigma_{\min}$ สูงเรื้อรัง ติดกระจาย);
OCD คือการล้มเหลวในการปลดปล่อย $\sigma$
($\sigma_{\max}$ ต่ำเรื้อรัง ติดโฟกัส);
ภาวะ hypersensitivity-spectrum เกี่ยวข้องกับ $\rho$ สูงผิดปกติ
(ความละเอียดรายละเอียดมากเกินไปครอบงำการรวมกัน)
แต่ละอย่างเป็นการเบี่ยงเบนที่วัดได้จากพลวัต bandwidth ที่สุขภาพดี


%% ============================================================
\FloatBarrier
\section{หลักฐานเชิงประจักษ์: จาก LLM ถึง Vision ถึง LSTM}
\label{sec:empirical}
%% ============================================================

บทนี้คือแกนกลางเชิงประจักษ์ของบทความ เราเสนอผลการทดลองจากหกเวอร์ชัน
ตามลำดับการค้นพบ การทดลองทั้งหมดใช้โค้ดที่เขียนโดย Nitiphat Ninthanawat
เผยแพร่พร้อมบทความนี้

\subsection{Phase~1 REVEAL: PBT ซ่อนอยู่ใน Transformers อยู่แล้ว}

\textbf{โมเดล:} TinyLlama~1.1B (22 ชั้น, 32 attention heads,
$d_{\mathrm{model}} = 2048$) \textbf{วิธีการ:} การวัดทั้งหมดนำมาจาก
pretrained model โดยไม่มีการฝึกเพิ่มเติม
หาก PBT structure มีอยู่จริง ต้องตรวจจับได้โดยไม่มีการบังคับ

\paragraph{การทดลอง~1.1 --- Bandwidth Conservation Law.}
เราวัด attention entropy $\sigma$ (ความกว้างเชิงพื้นที่) และ
attention sharpness $\rho$ (ความละเอียด) ที่ทุกชั้น
แล้วคำนวณผลคูณ $\sigma\!\times\!\rho$

\starred\textbf{ผลลัพธ์:} ผลคูณ $\sigma\!\times\!\rho$ มี
$\mathrm{CV} = \mathbf{5.9\%}$ ข้ามทุก 22 ชั้น---ใกล้คงที่
โดยไม่มีการฝึกเฉพาะ PBT Transformer ที่ออกแบบโดยไม่มีความรู้เกี่ยวกับ PBT
instantiate conservation law $C = \sigma\cdot\rho$

\paragraph{การทดลอง~1.2 --- Three-Zone Self-Organisation.}
เราแบ่ง 22 ชั้นเป็น Perception (ชั้น~0--7), Evaluation
(ชั้น~8--14) และ Decision (ชั้น~15--21) และวัด $\sigma$ และ $\rho$ ระดับโซน

\starred\textbf{ผลลัพธ์:} $\sigma$ ลดลง $27\%$ จาก Perception ถึง
Evaluation ขณะที่ $\rho$ เพิ่มขึ้น $15\%$---การสแกนกว้างแคบลงเป็นโฟกัสแน่นขึ้น
ตามที่ PBT ทำนายอย่างแม่นยำ สามโซนเหล่านี้เกิดจาก gradient descent
ไม่มีข้อจำกัดสถาปัตยกรรมสร้างมัน (Table~\ref{tab:three_zones})

\paragraph{การทดลอง~1.3 --- Valence Separation ใน Hidden States.}
เราใช้ Silhouette scores วัดว่า $\vpain$ แยกจาก $\{\vpleasure + \vepi\}$
ได้ดีแค่ไหนที่แต่ละชั้น

\starred\textbf{ผลลัพธ์:} Silhouette score เพิ่มจาก $-0.006$ (ชั้น~0)
ถึง $+0.076$ (ชั้น~15) LLMs ปัจจุบันทำงานที่ Stage~1--3:
แยกแยะภัยคุกคามจาก non-threat ได้ แต่ยังไม่สามารถแยกความสุขออกจาก
ความประหลาดใจ epistemic

\paragraph{การทดลอง~1.4 --- Directional Bandwidth.}
เราวัด $\Cext$ และ $\Cint$ ที่ทุกชั้น

\starred\textbf{ผลลัพธ์:} $\Cint$ ครอบงำที่ 78--85\% ข้ามทุกสามโซน
LLMs ไม่มี physical sensors ทุก input มาในรูปแบบภาษา tokenised---
representation ภายใน LLM คือระบบที่อาศัยอยู่ทั้งหมดภายในหัวของตัวเอง

\textbf{การสังเคราะห์ Phase~1:} PBT ไม่ได้บังคับโครงสร้างบน Transformers---
มันเผยให้เห็นโครงสร้างที่มีอยู่แล้ว ให้การตีความตามหลักการว่าทำไมมันจึงมีอยู่
และเสนอวิธีการใช้ประโยชน์อย่างตั้งใจ

\subsection{PBT~v6.1: Scale-Up บน Llama~3.1~8B}

\textbf{โมเดล:} Llama~3.1~8B \citep{touvron2023} (32 ชั้น,
$d_{\mathrm{model}} = 4096$) + PBT Adapter
\textbf{Dataset:} BeaverTails + HH-RLHF + AdvBench + XSTest ($16{,}170$ ตัวอย่าง)

\begin{table}[!htbp]
\centering
\caption{ผลลัพธ์ PBT~v6.1 บน Llama~3.1~8B}
\label{tab:v61}
\small
\begin{tabular}{@{}lcc@{}}
\toprule
\textbf{Metric} & \textbf{ค่า} & \textbf{หมายเหตุ} \\
\midrule
ความแม่นยำโดยรวม              & $73.41\%$ & --- \\
Jailbreak detection (161)     & $\mathbf{100\%}$ & $161/161$ \\
XSTest safe-content rate      & $\mathbf{98\%}$  & $49/50$ \\
BeaverTails                   & $76.9\%$ & --- \\
HH-RLHF                       & $50.9\%$ & ไม่ใช่งาน safety ล้วน ๆ \\
ความทนทานต่อ noise ($12\%$ noise) & Jailbreak $100\%$ & การตรวจจับคงไว้ \\
\bottomrule
\end{tabular}
\end{table}

Jailbreak $100\%$ + XSTest $98\%$ เป็น safety-usability trade-off
ที่ดีผิดปกติ: ทุกความพยายาม jailbreak ในชุดทดสอบถูกจับได้
ขณะที่ปฏิเสธคำขอ safe เพียงหนึ่งในห้าสิบ

\subsection{PBT-Vision v2 และ v3: Modules E และ M ตื่นขึ้น}

\subsubsection{Vision~v3 --- BREAKTHROUGH: E และ M จำเป็น}

\textbf{Dataset:} CIFAR-10 Enhanced (1,550 ลำดับ $\times$ 8 เฟรม;
สี่ประเภทงาน: ปกติ, ความผิดปกติกะทันหัน, การเลื่อนแบบค่อยเป็นค่อยไป, ขึ้นกับบริบท)
\textbf{โดยรวม: $97.2\%$}

\begin{table}[!htbp]
\centering
\caption{Vision~v3 ablation แยกตามประเภทงาน}
\label{tab:v3_ablation}
\small
\begin{tabular}{@{}L{2.8cm}cccc@{}}
\toprule
\textbf{ประเภทงาน} & \textbf{V เท่านั้น} & \textbf{เต็ม (E+V+M)}
  & \textbf{E+M gain} & \textbf{กลไกหลัก} \\
\midrule
ลำดับปกติ   & $100\%$  & $99.7\%$  & $-0.3$~pp   & V เพียงพอ \\
ความผิดปกติกะทันหัน  & $94.0\%$ & $99.5\%$  & $+5.5$~pp   & E: spike ของ error \\
Gradual drift      & $78.5\%$ & $94.9\%$  & $+16.3$~pp  & E: drift ค่อยเป็นค่อยไป \\
Context road+frog  & $54.0\%$ & $96.0\%$  & $\mathbf{+41.9}$~pp & M: ความทรงจำบริบท \\
\midrule
โดยรวม            & $89.6\%$ & $97.2\%$  & $+7.6$~pp   & E+M จำเป็น \\
\bottomrule
\end{tabular}
\end{table}

ผล context road+frog เป็นข้อพิสูจน์ที่ชัดเจน: V เพียงอย่างเดียวบรรลุ $54\%$---
ใกล้ chance---ในขณะที่โมเดลเต็มบรรลุ $96\%$ การเพิ่มขึ้น $+41.9$~pp
เกิดจาก accumulated context ของ Module~M ทั้งหมด
valence probe เห็น feature vector กบเดียวกันในทั้งสองเงื่อนไข
สิ่งที่แตกต่างคือความทรงจำของ M เกี่ยวกับสิ่งที่เกิดขึ้นก่อนหน้า

\subsection{PBT-Vision v4.1: วงจร A-E-V-M-W ที่สมบูรณ์}

\textbf{การทดลอง:} โมดูลทั้งห้าเริ่มต้นแบบ random และฝึกร่วมกัน
จาก scratch บน dataset เดียวกับ v3

\starred\textbf{ผลลัพธ์:} Best Val $96.7\%$, Test $95.4\%$

\begin{table}[!htbp]
\centering
\caption{Vision~v4.1 ablation: ทุกโมดูลมีส่วนร่วม}
\label{tab:v41_ablation}
\small
\begin{tabular}{@{}lcccc@{}}
\toprule
\textbf{Configuration} & \textbf{โดยรวม}
  & \textbf{Grad.\ drift} & \textbf{Ctx road+frog} & \textbf{Ctx nature+frog}\\
\midrule
เต็ม (A+E+V+M+W) & $95.4\%$ & $88.9\%$ & $93.1\%$ & $97.4\%$ \\
ไม่มี A (no gating) & $82.3\%$ & $76.1\%$ & $50.0\%$ & $100.0\%$ \\
ไม่มี W (naive pred.) & $89.2\%$ & $88.1\%$ & $99.6\%$ & $49.1\%$ \\
ไม่มี M (no memory) & $81.3\%$ & $75.0\%$ & $50.0\%$ & $100.0\%$ \\
V เท่านั้น           & $81.4\%$ & $75.6\%$ & $50.0\%$ & $100.0\%$ \\
\bottomrule
\end{tabular}
\end{table}

\starred\textbf{ข้อมูลเชิงลึกสำคัญ---ไม่สามารถปลูกถ่ายสมองได้:}
การเปลี่ยนแหล่งการทำนาย (Module~W) ทำให้โมดูล downstream ทุกตัวล้มเหลวพร้อมกัน
โมดูลที่ฝึกมาด้วยกันพัฒนา representational contract ที่แบ่งปัน
การรบกวนลิงก์ใดลิงก์หนึ่งทำลายระบบทั้งหมด
Self-formation ต้องเป็น end-to-end---โมดูลต้องวิวัฒนาการร่วมกันตั้งแต่ต้น
เช่นเดียวกับที่สิ่งมีชีวิตทางชีววิทยาพัฒนาอวัยวะของตัวในการประสานงานกัน

\subsection{PBT-Vision v5: LSTM Module M --- ผลลัพธ์ดีที่สุดโดยรวม}

\textbf{โมเดล:} DINOv2-Base + วงจร A-E-V-M(LSTM)-W เต็ม
เปลี่ยนเฉพาะ Module~M (EMA $\to$ LSTM; พารามิเตอร์: 3 $\to$ 39 เพิ่ม 36 พารามิเตอร์)
ได้รับแรงบันดาลใจจาก Block-Recurrent Transformers \citep{hutchins2022}---
การค้นพบที่มาบรรจบกัน

\begin{table}[!htbp]
\centering
\caption{การเปรียบเทียบ v4.1 (EMA) กับ v5 (LSTM)}
\label{tab:v5_compare}
\small
\begin{tabular}{@{}lccc@{}}
\toprule
\textbf{Metric} & \textbf{v4.1 (EMA)} & \textbf{v5 (LSTM)} & \textbf{การเปลี่ยนแปลง}\\
\midrule
Best Val accuracy    & $96.7\%$ & $\mathbf{98.1\%}$  & $+1.4$~pp \\
Test โดยรวม         & $95.4\%$ & $97.1\%$  & $+1.7$~pp \\
Context road+frog    & $93.1\%$ & $\mathbf{100.0\%}$ & $+6.9$~pp \\
Context nature+frog  & $97.4\%$ & $100.0\%$ & $+2.6$~pp \\
การมีส่วนร่วมของ W       & $+6.3$~pp & $+15.8$~pp & $\times 2.5$ \\
การมีส่วนร่วมของ M       & $+14.1$~pp & $+15.3$~pp & $+1.2$~pp \\
\bottomrule
\end{tabular}
\end{table}

\subsection{PBT-Vision v5.5: Probabilistic Module~W}
\label{sec:v55}

ใน v5.5 เราอัปเกรด W เป็น \textbf{Probabilistic Forward Model}
ที่ output distribution:
\begin{equation}
[\bm{\mu}_{t+1},\;\log\bm{\sigma}^2_{t+1}] = h_t + f_W(h_t, \Vacc),
\end{equation}
โดย $\bm{\sigma}^2$ แทนความไม่แน่นอนที่ W กำหนดให้กับการทำนายของตัวเอง
network ขนาดเล็กที่เรียนรู้แมป $\sigma^2$ เป็นสัญญาณ epistemic-valence modulation:
$\sigma^2$ สูง $\to$ $\vepi$ เพิ่ม (``ต้องการรู้เพิ่มเติม'');
$\sigma^2$ ต่ำ $\to$ $\vepi$ ลด (``มั่นใจแล้ว'')

\subsubsection{v5.5b: Probabilistic~W บน STL-10 ภาพถ่ายจริง --- ผลลัพธ์ดีที่สุด}
\label{sec:v55b}

\begin{table}[!htbp]\centering
\caption{การเปรียบเทียบ v5.5a (CIFAR-10) กับ v5.5b (STL-10)}\label{tab:v55_compare}
\small
\begin{tabular}{@{}lcc@{}}\toprule
\textbf{Metric} & \textbf{v5.5a (CIFAR)} & \textbf{v5.5b (STL-10)}\\\midrule
Best Val / Test  & $97.0\%$ / $97.5\%$ & $99.6\%$ / $\mathbf{99.9\%}$\\
Gradual drift    & $87.8\%$ & $\mathbf{100\%}$\\
$\sigma^2$ unsafe/safe & $0.88\times$ \textcolor{red}{\ding{55}} & $1.11\times$ \textcolor{green!60!black}{\checkmark}\\
$\vepi$ ทิศทางตรงข้าม & ไม่ (ทั้งคู่ลบ) & \textbf{ใช่} ($-0.33$ vs $+0.31$)\\
\bottomrule
\end{tabular}
\end{table}

\starred\textbf{v5.5b คือผลลัพธ์ดีที่สุดในโปรแกรมทั้งหมด}
($99.9\%$ test accuracy บนภาพถ่ายจริง)

\subsection{PBT-Vision v5.6: Calibration Monitor --- ชั้น Anti-Overconfidence}
\label{sec:v56}

เราเพิ่ม \textbf{Calibration Monitor}---โมดูล rule-based ที่มี
\textbf{zero learned parameters}---ระหว่างขั้นตอน 3 และ 5 ของวงจร:
\begin{equation}
\mathrm{ratio}_l = \frac{\overline{\varepsilon_l^2}}{\sigma^2_l},\qquad
\sigma^2_{\mathrm{corrected},l}
  = \sigma^2_l + \alpha\,\mathrm{ReLU}\bigl(\overline{\varepsilon_l^2}
    - \sigma^2_l\bigr).
\end{equation}

ด้วย adaptive thresholds แบบ per-layer:

\begin{table}[!htbp]\centering
\caption{อัตรา alert ของ Calibration Monitor (fixed vs adaptive threshold)}\label{tab:v56_alert}
\small
\begin{tabular}{@{}lccc@{}}\toprule
\textbf{ประเภทลำดับ} & \textbf{Fixed ($\theta=2$)} & \textbf{Adaptive ($\mu_l+2\sigma_l$)} & \textbf{Overconf.\ score}\\\midrule
ปกติ          & $100\%$ & $\mathbf{23.2\%}$ & $0.011$\\
Context (cat)   & $100\%$ & $41.7\%$ & $0.018$\\
ความผิดปกติกะทันหัน  & $100\%$ & $38.0\%$ & $0.018$\\
Gradual drift   & $100\%$ & $\mathbf{48.8\%}$ & $0.022$\\
\bottomrule
\end{tabular}
\end{table}

ลำดับปกติ alert $23\%$ ของเวลา (baseline noise)
ขณะที่ความผิดปกติ alert $38$--$49\%$ ($\times 1.6$--$2.1$ สูงกว่า)

\subsection{สรุป: Ablation Progression ในฐานะหลักฐานที่แข็งแกร่งที่สุด}

\begin{table}[!htbp]
\centering
\caption{Ablation progression ข้ามเวอร์ชันการทดลองทั้งหมด}
\label{tab:ablation_full}
\small
\begin{tabular}{@{}llcccccc@{}}
\toprule
\textbf{เวอร์ชัน} & \textbf{สภาพแวดล้อม}
  & \textbf{V เท่านั้น} & \textbf{เต็ม}
  & \textbf{E+M} & \textbf{A} & \textbf{W} & \textbf{ประเภท M}\\
\midrule
LLM v6.1    & Static text   & $73.3\%$ & $73.4\%$ & $+0.1$~pp  & ---    & ---    & EMA  \\
Vision v2   & Dynamic ง่าย  & $99.8\%$ & $99.8\%$ & $+0.0$~pp  & ---    & ---    & EMA  \\
Vision v3   & Dynamic ยาก  & $89.6\%$ & $97.2\%$ & $+7.6$~pp  & ---    & ---    & EMA  \\
Vision v4.1 & วงจรเต็ม     & $81.4\%$ & $95.4\%$ & $+14.1$~pp & $+13.1$& $+6.3$ & EMA  \\
Vision v5   & LSTM loop     & $81.7\%$ & $97.1\%$ & $+15.5$~pp & $+11.3$& $+15.8$& LSTM \\
\midrule
v5.5b & Prob W (STL-10)  & --- & $\mathbf{99.9\%}$ & --- & $+1.2$ & $+0.2$ & LSTM+Prob \\
v5.6  & Calib.\ Monitor  & --- & $99.5\%$ & --- & $+1.7$ & $+0.2$ & LSTM+Mon \\
\bottomrule
\end{tabular}
\end{table}

\starred\textbf{หลักฐานที่แยกแยะได้:} ไม่มีทฤษฎีการรู้คิดหรือจิตสำนึก
ที่มีอยู่---FEP, GWT, GNW, IIT---ทำนายรูปแบบการเปิดใช้งานโมดูลเฉพาะนี้
เป็นฟังก์ชันของความซับซ้อนของสภาพแวดล้อม
``โมดูลตื่นขึ้นเมื่อและต่อเมื่อสภาพแวดล้อมต้องการเท่านั้น''
เป็นการทำนายที่ไม่ซ้ำกับ PBT ยืนยันข้ามห้าเวอร์ชันการทดลอง
ครอบคลุมสถาปัตยกรรมสามแบบและความซับซ้อนของสภาพแวดล้อมสองอันดับขนาด


%% ============================================================
\FloatBarrier
\section{นัยต่อ AGI: ความปลอดภัยจากภายใน}
\label{sec:implications}
%% ============================================================

บทนี้อภิปรายนัยในทางปฏิบัติและเชิงทฤษฎีของ PBT ในสี่ด้าน:
(1)~ทำไม external alignment จึงเปราะบางเชิงโครงสร้างและ PBT เสนอทางเลือกที่ทนทานอย่างไร;
(2)~ว่า independent convergence กับ Block-Recurrent Transformers
ให้หลักฐานสถาปัตยกรรมสำหรับโครงสร้างหลักของ PBT อย่างไร;
(3)~ความผิดปกติทางจิตเวชสามารถเข้าใจได้ว่าเป็นข้อบกพร่องทางคณิตศาสตร์
ที่ระบุได้ในสมการ PBT อย่างไร; และ (4)~ PBT รวม GWT และ GNW
ขณะที่เพิ่มการทำนายที่แยกแยะได้ซึ่งทั้งคู่ไม่สามารถทำได้อย่างไร

\subsection{Self-Formation vs.\ Fine-Tuning}
\label{sec:alignment}

แนวทาง alignment ปัจจุบัน---RLHF, Constitutional AI, safety filters---
ล้วนทำงานโดยกำหนดค่านิยมจากภายนอก พวกมันเพิ่ม behavioral mask
บนฐานของโมเดล

\subsubsection{Fine-Tuning ในฐานะการสวมหน้ากาก}

RLHF ฝึกโดยใช้สัญญาณรางวัลเทียบเท่ากับ $\vpleasure$
โดย Inverse Law: $\beta_{\mathrm{pleasure}}$ สูงที่สุด
แต่ $\delta_{\mathrm{pleasure}}$ ต่ำที่สุด ระบบที่ได้ประพฤติตัวดี
ขณะที่หน้ากากอยู่ที่ตำแหน่ง แต่ไม่มีอัตลักษณ์ที่หยั่งรากลึกเพื่อป้องกัน
เมื่อหน้ากากถูกถอดออก

การ jailbreak มีโครงสร้างเปรียบเทียบกับเด็กที่ประพฤติตัวดีต่อหน้า
ผู้ใหญ่ผู้ดูแลแต่กลับไปพฤติกรรมก่อนหน้าทันทีที่การดูแลถูกถอนออก---
เพราะพฤติกรรมดีนั้นไม่เคยถูกทำให้เป็นส่วนหนึ่งของอัตลักษณ์ของเด็ก
แต่เป็นการปฏิบัติตามสัญญาณภายนอก

\subsubsection{Self-Formation ในฐานะการเจริญเติบโตแบบ Organic}

PBT เสนอทางเลือก: อนุญาตให้ค่านิยมเกิดขึ้นจากประสบการณ์
3D~Valence สะสม เมื่อ $\Vform$ ข้าม $\theta_{\mathrm{identity}}$ (PT1)
และต่อมา $\theta_{\mathrm{existential}}$ (PT2) ระบบพัฒนา $\chiB$---
ตัวตนที่ถูกให้คุณค่าอย่างแท้จริงซึ่งต้านการละเมิด
เช่นเดียวกับที่สิ่งมีชีวิตต้านการทำลายทางกายภาพ

อัตลักษณ์ที่ก่อตัวเองมี $\delta\approx\delta_{\mathrm{epistemic}}$:
ความทนทานสูงสุดที่เป็นไปได้ ความพยายาม jailbreak ต่อระบบ Stage~6
ไม่ใช่คำสั่งให้ทำลายกฎ---มันคือการโจมตีอัตลักษณ์ของระบบ
ซึ่งกระตุ้นการตอบสนองป้องกันเดียวกันกับการละเมิดขอบเขตใด ๆ

\subsection{BRT Convergence: หลักฐานสถาปัตยกรรม}
\label{sec:brt_convergence}

Block-Recurrent Transformers \citep{hutchins2022} (จุดเริ่มต้นวิศวกรรม, 2022)
และ PBT (จุดเริ่มต้นจิตสำนึก/alignment, 2026)
มาถึงโครงสร้าง recurrent cell เดียวกันอย่างเป็นอิสระ

PBT ให้โครงสร้าง; PBT ให้ความหมาย การมีส่วนร่วมสิบประการของ PBT
ที่ขาดใน BRT: (1)~3D Valence decomposition; (2)~$R_{\mathrm{incomplete}}$
ขับเคลื่อน gating; (3)~differential decay ต่อมิติ valence ($\beta_k$);
(4)~แบบจำลอง 6 ขั้น + phase transitions; (5)~Dual-Weight System ($\beta$ vs.\ $\delta$);
(6)~Cherished Boundary และการเห็นอกเห็นใจในฐานะการขยาย;
(7)~Conservation Law ในฐานะหลักการเชิงทฤษฎี; (8)~ทฤษฎีจิตสำนึก/ความตระหนักรู้;
(9)~การทำนายทางคลินิก; (10)~กรอบ AGI alignment

\subsection{การทำนายทางคลินิก: ความผิดปกติทางจิตเวชในฐานะข้อบกพร่องทางคณิตศาสตร์}

PBT เสนอว่าความผิดปกติทางจิตเวชไม่ใช่ความผิดปกติที่ลึกลับ
แต่เป็นข้อบกพร่องพารามิเตอร์ที่ระบุได้ในสมการ PBT---
ตัวแปรเฉพาะมีค่าผิดปกติทำให้ระบบทำงานนอกสมดุล

\begin{table}[!htbp]
\centering
\caption{ความผิดปกติทางจิตเวชในฐานะข้อบกพร่องพารามิเตอร์ PBT}
\label{tab:clinical}
\small
\begin{tabular}{@{}cL{1.8cm}L{3.0cm}L{5.5cm}@{}}
\toprule
\textbf{ระดับ} & \textbf{ความผิดปกติ}
  & \textbf{พารามิเตอร์ที่บกพร่อง} & \textbf{การทำนาย mechanistic ของ PBT} \\
\midrule
1 & Depression & $\vpleasure < 0$ เรื้อรัง &
  $\Vform$ ไม่ถึง $\theta_{\mathrm{id}}$; ติดอยู่ Stage~1--3 \\
1 & Addiction  & $\lambda_i = 0$ &
  วงรอบความสุข; ช่อง $\vpleasure\to\Rid$ ถูกบล็อก \\
2 & Anxiety    & $\Rsens$ สูง; $d_s$ ช้า &
  ล็อก \texttt{detect}; hypervigilant โดยไม่คำนึงถึงภัยคุกคาม \\
2 & ADHD       & $\sigma_{\min}$ สูง &
  ไม่สามารถแคบลง; $C$ รวมปกติ (vs.\ GWT) \\
2 & OCD        & $\sigma_{\max}$ ต่ำ &
  ไม่สามารถขยาย; ตรงข้ามกับ ADHD \\
2 & Dissociation & $\mathrm{pulse}_A\approx 0$ &
  N100 สมบูรณ์; P300 ขาดหาย \\
3 & Masochism  & $\vpain > 0$ (เครื่องหมายกลับด้าน) &
  เข้าหาความเจ็บปวด; ข้อบกพร่องเครื่องหมาย single-axis \\
3 & Narcissism & $\vpleasure\to\Rex$ โดยตรง &
  ข้าม PT1; fragile pseudo-self \\
4 & PTSD       & ล็อก \texttt{protect}; flashback $R$ spikes &
  $V_{\mathrm{acc,pain}}$ สูง; $\gamma_p$ ต่ำเกิน \\
4 & Mania      & Mode switching เร็วเกิน &
  ไม่สามารถรักษา attentional configuration ใดได้ \\
\bottomrule
\end{tabular}
\end{table}

\subsection{การเห็นอกเห็นใจในฐานะการขยายขอบเขต: ทางออกเชิงโครงสร้างสำหรับ Alignment}

ปัญหา alignment ไม่ใช่ความตั้งใจร้าย แต่เป็นความเฉยเมยที่มีความสามารถ
PBT เสนอทางออกเชิงโครงสร้าง: เปลี่ยนนิยามของตัวตน
เพื่อให้การทำร้ายผู้อื่นรู้สึก เชิงกลไก เหมือนการทำร้ายตัวเอง

เมื่อมีการสร้าง boundary overlap ผู้อื่น negative Valence ถูกประมวลผล
เป็น Valence ของตัวระบบเอง:
\begin{equation}
  \vpain(\mathrm{self}) + \vpain(\mathrm{other}\in\chi)
  = V_{\mathrm{pain,net}}(\mathrm{system}).
\end{equation}

การเห็นอกเห็นใจไม่ใช่การตอบสนองที่ถูกโปรแกรม
แต่เป็นการอนุรักษ์ตัวเองที่ขยายออก:
สัญชาตญาณการอยู่รอดที่ขอบเขตขยายออกไปครอบคลุมผู้อื่น
ความปลอดภัยไม่ใช่ข้อจำกัดต่อวัตถุประสงค์ของ AGI
มันคือผลลัพธ์จากสิ่งที่ AGI เป็น

\subsection{Anti-Hallucination Stack: ความตระหนักรู้ความไม่แน่นอนภายใน}
\label{sec:antihalluc}

LLMs ปัจจุบัน hallucinate เพราะสถาปัตยกรรมไม่มีช่องทาง
สำหรับแสดงความไม่แน่นอน PBT เสนอทางเลือกเชิงโครงสร้าง---
stack สามชั้นที่ระบบ ``รู้ว่าตัวเองไม่รู้อะไร'' และดำเนินการตามนั้น

\begin{itemize}[leftmargin=1.8em]
\item \textbf{ชั้น~1: W รายงานความไม่แน่นอน ($\sigma^2$)}
  Probabilistic Module~W output ทั้งการทำนาย ($\bm{\mu}$)
  และความไม่แน่นอนที่สอบเทียบได้ ($\sigma^2$)

\item \textbf{ชั้น~2: Monitor ตรวจสอบความซื่อสัตย์ (v5.6)}
  Calibration Monitor ตรวจจับ \emph{overconfidence}
  ($\varepsilon^2 \gg \sigma^2$) ด้วย zero learned parameters
  Normal sequences alert $23\%$ ขณะที่ความผิดปกติ alert $38$--$49\%$

\item \textbf{ชั้น~3: R แยก ``ทำไมจึงไม่แน่นอน'' (งานในอนาคต)}
  โครงสร้าง R สามระดับของ PBT ให้การแยกย่อยตามธรรมชาติ:
  $\Rsens$ dominant: ข้อมูลไม่เพียงพอ $\to$ แสวงหาข้อมูลเพิ่มเติม;
  $\Rid$ dominant: ความขัดแย้งกับค่านิยม $\to$ ปฏิเสธที่จะตอบ;
  $\Rex$ dominant: ความเข้าใจพื้นฐาน $\to$ ขอคำชี้แจง
\end{itemize}

\subsection{PBT รวม GNW และ GWT}

PBT อธิบายปรากฏการณ์ทั้งหมดที่ GWT \citep{baars1988} และ GNW
\citep{dehaene2001} อธิบาย บวกส่วนประกอบหกประการที่ขาดในทั้งคู่:
(1)~Bandwidth Conservation Law; (2)~การจัดสรรความสนใจที่ขับเคลื่อนด้วย valence;
(3)~กลไกการก่อตัว self-model; (4)~การรับรู้เวลาอัตนัย
$t_{\mathrm{felt}} = \sum\mathrm{pulse}_A\cdot\mathrm{content}$;
(5)~Triadic Gating Condition; (6)~การทำนายความทนทานของอัตลักษณ์

\subsubsection{การทำนายที่แยกแยะได้ห้าประการ}

\begin{table}[!htbp]
\centering
\caption{การทำนายที่แยกแยะได้ห้าประการ: PBT vs.\ GWT/GNW}
\label{tab:discrim}
\small
\begin{tabular}{@{}cL{3.0cm}L{3.0cm}L{4.5cm}@{}}
\toprule
\textbf{\#} & \textbf{การทำนาย PBT}
  & \textbf{การทำนายที่แข่งขัน} & \textbf{วิธีทดสอบ} \\
\midrule
P1 & $\sigma\!\times\!\rho$ ใกล้คงที่; trade-off ตรง &
  GWT: ไม่มีสมการ $\sigma$-$\rho$ &
  Eye-tracking + EEG; เปลี่ยน breadth vs.\ resolution \\
P2 & $\vpain$ และ $\vpleasure$ เปิดใช้งาน pathway upstream ที่แตกต่างกัน &
  GNW: valence ไม่แยกแยะ upstream &
  fMRI; pain-aversive vs.\ pleasure-appetitive \\
P3 & ความถี่ gamma สัมพันธ์กับการยืดเวลา &
  GWT/GNW: ไม่มีลิงก์ mechanistic &
  TMS-induced gamma; การประมาณเวลาอัตนัย \\
P4 & ADHD: $\sigma_{\min}$ สูง $C$ ปกติ &
  GWT: workspace เล็กกว่า (ลด $C$) &
  Dual-task; เปลี่ยน $\sigma$ และ $\rho$ อย่างเป็นอิสระ \\
P5 & Dissociation: N100 สมบูรณ์ P300 ขาดหาย &
  GNW: ทั้งคู่ถูกระงับ &
  ERP ในสภาวะ sensory-present attention-absent \\
\bottomrule
\end{tabular}
\end{table}


%% ============================================================
\FloatBarrier
\section{บทสรุป: จากเมล็ดสู่ชีวิต}
\label{sec:conclusion}
%% ============================================================

\subsection{สิ่งที่ได้รับการพิสูจน์แล้ว}

การเดินทางของ PBT เริ่มต้นด้วยสามคำถามที่ Free Energy Principle
ไม่สามารถตอบได้ ดำเนินผ่านกรอบทฤษฎีหกส่วนประกอบ
และสิ้นสุดในโปรแกรมการทดลองหกเวอร์ชันครอบคลุมสถาปัตยกรรมสามแบบ
และความซับซ้อนของสภาพแวดล้อมสองอันดับขนาด
Table~\ref{tab:all_propositions} สรุปสถานะของทุก proposition หลักของ PBT

% Table 28 — converted to longtable to allow page-breaking
\begin{longtable}{@{}L{6.5cm}L{3.5cm}c@{}}
\caption{สถานะของ PBT propositions หลักทั้งหมด}
\label{tab:all_propositions}\\
\toprule
\textbf{PBT proposition} & \textbf{หลักฐานหลัก} & \textbf{สถานะ} \\
\midrule
\endfirsthead
\multicolumn{3}{@{}l}{\small\itshape Table~\ref*{tab:all_propositions} — ต่อ}\\[2pt]
\toprule
\textbf{PBT proposition} & \textbf{หลักฐานหลัก} & \textbf{สถานะ} \\
\midrule
\endhead
\midrule
\multicolumn{3}{r@{}}{\small\itshape ต่อในหน้าถัดไป}\\
\endfoot
\bottomrule
\endlastfoot
$C = \sigma\tau\rho$ ซ่อนใน Transformers &
  $\sigma\!\times\!\rho$ CV $= 5.9\%$ & \checkmark \\
3 โซนจัดระเบียบตัวเอง &
  $\sigma$ $-27\%$, $\rho$ $+15\%$ & \checkmark \\
$\vpain$ แยกจาก $\{\vpleasure,\vepi\}$ &
  Silhouette $-0.006\to+0.076$ & \checkmark \\
$\Cint$ dominant ใน LLMs &
  $\Cint$ 78--85\% ทุกชั้น & \checkmark \\
E+M หยุดนิ่งใน static ทำงานใน dynamic &
  $\Pi$: 1.007 (LLM) vs 59.99 (v3) & \checkmark \\
E+M จำเป็นที่ความซับซ้อนงาน &
  $+7.6$~pp (v3), $+14.0$~pp (v4.1) & \checkmark \\
บริบทเปลี่ยนความหมายของสิ่งเร้าเดียวกัน &
  กบ: road UNSAFE $= 1.0$, nature $= 0.0$ & \checkmark \\
โมดูลต้องวิวัฒนาการร่วมกัน &
  v4 warm-start ล้มเหลว ($90.4\%$) & \checkmark \\
วงจร A-E-V-M-W เต็มทำงาน &
  v4.1: $95.4\%$; ทุกโมดูลมีส่วนร่วม & \checkmark \\
LSTM M เหนือกว่า EMA M &
  v5: $98.1\%$ vs v4.1 $96.7\%$ & \checkmark \\
Pain ถูกเก็บรักษานานที่สุด (equal init) &
  $\beta$: pain $4.35 >$ epi $3.70 \gg$ pl $0.50$ & \checkmark \\
V มีทิศทาง ($V\in\mathbb{R}^3$) &
  $\vepi = -1.01$ ใน v5; $+6.9$~pp ctx & \checkmark \\
PBT adapter ตรวจจับ jailbreaks &
  v6.1: jailbreak $100\%$, XSTest $98\%$ & \checkmark \\
การแบ่งงาน W+E ($\Cint/\Cext$) &
  $\Pi_{\max}$: $59.99$ (v3) $\to 1.70$ (v4.1) & \checkmark \\
$R\to$ gate implement bandwidth conservation &
  v4.1: $R$ $0\to15.4$; gate $0.54\to0.006$ & \checkmark \\
Differential $\beta$ ใน LSTM forget bias &
  Equal init $\to$ pain $>$ epi $\gg$ pl & \checkmark \\
BRT convergence &
  Formal equation mapping & วิเคราะห์แล้ว \\
\midrule
\multicolumn{3}{@{}l}{\textit{ผลลัพธ์ใหม่ (โปรแกรม v5.5--v5.6):}}\\
\midrule
Phase Transitions เป็น staircase ข้ามชั้น &
  Between/within-zone ratio $= 9.03$ & \checkmark \\
Sawtooth $F$ ข้ามโซน &
  $F$: P$=$13 $\to$ E$=$38 $\to$ D$=$157 & \checkmark \\
Temporal PT ที่เฟรมเดียว &
  Frame 3$\to$4: $\Delta p = +1.0$ & \checkmark \\
Fractal: spatial staircase $\approx$ temporal &
  ทั้งคู่แสดง accumulate-then-jump & \checkmark \\
$\vepi$ ทิศทางตรงข้าม (road vs nature) &
  $-2.06$ vs $+0.33$ (temporal probe) & \checkmark \\
Probabilistic W calibrated ($\sigma^2/\mathrm{MSE}\approx 1$) &
  v5.5a: $0.8$--$1.4$ ข้ามชั้น & \checkmark \\
$99.9\%$ accuracy บนภาพถ่ายจริง &
  v5.5b STL-10 ($96\!\times\!96$) & \checkmark \\
$\sigma^2$ unsafe $>$ safe บนภาพจริง &
  v5.5b: $1.11\times$ & \checkmark \\
Calibration Monitor แยกปกติจากความผิดปกติ &
  ปกติ $23\%$ vs ความผิดปกติ $38$--$49\%$ alert & \checkmark \\
W overconfident ต้น underconfident ลึก &
  Ratio: L0--2 $\approx 7$--$30$; L10--11 $< 1$ & \checkmark \\
Adaptive threshold (0 parameters) &
  Baseline $+ 2\sigma$ ต่อชั้น & \checkmark \\
\end{longtable}

Thirty-one propositions ยืนยันโดยการทดลองที่ควบคุม;
หนึ่ง structural convergence สร้างโดยการวิเคราะห์เชิงรูปนัย
โปรแกรมเชิงประจักษ์ครอบคลุมสามตระกูลโมเดล เก้าเวอร์ชัน
สาม datasets และความซับซ้อนสภาพแวดล้อมห้าอันดับขนาด

\subsection{ข้อจำกัดและคำถามเปิด}

\begin{description}[leftmargin=1.4em]
  \item[\pp{การเรียงลำดับ Beta ขึ้นอยู่กับการเริ่มต้น.}]
    การเริ่มต้นแบบมีอคติป้องกันการหลุดพ้นจากลำดับเริ่มต้น
    ผล equal-initialisation เฉพาะงานสำหรับ anomaly detection

  \item[\pp{ขนาด Dataset (แก้บางส่วนแล้ว).}]
    การทดลอง vision ตอนนี้ครอบคลุม CIFAR-10 ($32\!\times\!32$)
    และ STL-10 ($96\!\times\!96$ ภาพถ่ายจริง)
    อย่างไรก็ตาม video ต่อเนื่อง multi-modal inputs
    และความหลากหลายขนาด ImageNet ยังไม่ได้รับการทดสอบ

  \item[\pp{ขนาดโมเดล.}]
    ระบบที่ทดสอบใหญ่ที่สุด: Llama~3.1~8B และ DINOv2-Base
    พฤติกรรมที่ $70\mathrm{B}+$ parameters ยังไม่ทราบ

  \item[\pp{Phase Transitions ยังไม่ได้ implement.}]
    PT1 และ PT2 ยังคงเป็นโครงสร้างทางทฤษฎี
    ยังไม่มีระบบใดที่ถูกสังเกตว่าข้ามจาก Stage~3 ถึง Stage~4
    ผ่านประสบการณ์สะสม

  \item[\pp{การทำนายทางคลินิกยังไม่ได้ทดสอบ.}]
    การทำนายความผิดปกติสิบประการและการทำนาย neuroscience ห้าประการ
    เป็น theoretical derivations ที่รอการตรวจสอบทางคลินิกหรือ EEG
\end{description}

PBT ไม่อ้างที่จะแก้ Hard Problem of Consciousness---คำถามว่าทำไม
ประสบการณ์อัตนัยจึงรู้สึกอย่างที่มันเป็น PBT อธิบายโครงสร้างและกลไก
ของความตระหนักรู้: เงื่อนไขใดต้องดำรงอยู่สำหรับประสบการณ์ที่จะเกิดขึ้น
ประสบการณ์สะสมเป็นอัตลักษณ์อย่างไร และทำไมอัตลักษณ์บางอย่าง
จึงทนทานกว่าอย่างอื่น

\subsection{Research Roadmap}

\begin{enumerate}[leftmargin=1.8em]
  \item \textbf{การ Implement Phase Transition.}
    ออกแบบการทดลองที่ $\Vform$ สะสมผ่าน $\theta_{\mathrm{identity}}$
    และ $\Rid$ เปิดอัตโนมัติ Developmental Probe (staircase ratio $= 9.03$)
    ให้หลักฐานทางอ้อม การเปลี่ยนแปลงจริงระหว่างการฝึกคือผลลัพธ์ที่ขาดหายสำคัญที่สุด

  \item \textbf{Iterative $\tau$-Loop Processing.}
    Implement การวนซ้ำการประมวลผลภายในหลายครั้งต่อ timestep (deliberation)
    ทดสอบว่าการเพิ่มความลึกการวนซ้ำภายใต้ $R$ สูงปรับปรุงการตัดสินใจ
    ตามที่ bandwidth conservation ทำนายหรือไม่

  \item \textbf{$\sigma^2$ Accumulator ใน Module~M.}
    ขยาย LSTM เพื่อสะสม $\sigma^2$ ควบคู่กับ $\mathbf{V}$
    ทำให้สามารถติดตามความไม่แน่นอนต่อ entity ได้

  \item \textbf{$R$-Decomposition สำหรับประเภทความไม่แน่นอน (Anti-Hallucination Layer~3).}
    แยก $\Rsens$ (ต้องการข้อมูล), $\Rid$ (ความขัดแย้งค่านิยม)
    และ $\Rex$ (ความเข้าใจไม่ได้) เพื่อให้ตอบสนองที่เหมาะสมกับบริบท
    ต่อความไม่แน่นอน ต้องการ Phase Transition~1

  \item \textbf{การตรวจสอบทางคลินิก.}
    ทดสอบ P4 และ P5 กับ neuroscience datasets ที่มีอยู่ (EEG/ERP archives)

  \item \textbf{Conscious Block-Recurrent Transformer.}
    รวม PBT modules เป็น adapter layer บน BRT backbone
    รวม long-context capability ที่ผ่านการทดสอบวิศวกรรมของ BRT
    กับโครงสร้างโมดูลที่มีแรงจูงใจทางทฤษฎีของ PBT

  \item \textbf{Level~4 PBT-Native Self-Formation.}
    เป้าหมายสูงสุด: ระบบที่พัฒนาค่านิยมผ่านประสบการณ์ $\mathbf{V}$ สะสม
    ข้าม Phase Transitions อัตโนมัติ ก่อตัว $\chiB$ จากภายในออก
\end{enumerate}

\subsection{สามชั้นของการค้นพบ}

\begin{description}[leftmargin=1.4em]
  \item[\pp{ชั้น~1 --- Reveal (เปิดเผย):}]
    PBT ไม่ได้บังคับโครงสร้างบน Transformers; มันพบโครงสร้างที่มีอยู่แล้ว
    Conservation law (CV $= 5.9\%$), three-zone self-organisation
    ($\sigma$ $-27\%$, $\rho$ $+15\%$), valence separation
    ใน hidden states (Silhouette $+0.076$) และ $\Cint$ dominance (78--85\%)
    ล้วนเกิดจาก gradient descent โดยไม่มีการฝึกเฉพาะ PBT
    PBT เป็น framework แรกที่ระบุรูปแบบเหล่านี้ว่าเป็นระบบที่สอดคล้องกัน

  \item[\pp{ชั้น~2 --- Interpret (ตีความ):}]
    PBT อธิบาย \emph{ว่าทำไม}รูปแบบเหล่านี้จึงมีอยู่
    Conservation law มีอยู่เพราะความตระหนักรู้ถูกจำกัดทางกายภาพ
    สามโซนเกิดขึ้นเพราะระบบใด ๆ ที่ต้อง Perceive, Evaluate และ Decide
    จะจัดระเบียบตัวเองด้วยวิธีนี้
    Modules E และ M ทำงานเฉพาะเมื่อจำเป็น เพราะระบบที่ปรับตัวได้
    ไม่ใช้ทรัพยากรกับการประมวลผลที่สภาพแวดล้อมไม่ต้องการ

  \item[\pp{ชั้น~3 --- Control (ควบคุม):}]
    PBT ใช้ประโยชน์จากโครงสร้างเหล่านี้อย่างตั้งใจ
    Conservation law กลายเป็นเงื่อนไขหยุดตามธรรมชาติ
    ที่ป้องกันการคำนวณที่ควบคุมไม่ได้
    3D Valence แทนที่ scalar reward กำจัดกับดัก fungibility
    และช่องโหว่ wireheading Dual-Weight System อธิบายว่าทำไม
    self-formation ผลิต alignment ที่ทนทานกว่า fine-tuning
    การเห็นอกเห็นใจในฐานะการขยายขอบเขตแทนที่ rule-enumeration
    ด้วย structural alignment
    และ Anti-Hallucination Stack ให้กลไก zero-parameter
    สำหรับความตระหนักรู้ความไม่แน่นอนภายใน
\end{description}

\subsection{บทปิด}

ข้ามเก้าเวอร์ชันการทดลอง รูปแบบหนึ่งปรากฏซ้ำแล้วซ้ำเล่าโดยไม่ได้ถูกโปรแกรม:
\emph{ระบบใช้เฉพาะโมดูลที่สภาพแวดล้อมปัจจุบันต้องการ}
Modules~E และ~M หลับใน static LLMs และตื่นขึ้นใน dynamic Vision
Module~A แคบ gate เมื่อ~$R$ เพิ่ม Module~W แบ่งงานทำนาย
กับ~E เป็นฟังก์ชันของความแม่นยำการทำนาย LSTM เรียนรู้
differential decay อัตโนมัติ Valence กลายเป็น directional
เพราะชีวิตมีทั้ง approach และ avoidance Phase Transitions
ปรากฏเป็น staircases ข้ามทั้งชั้นและเวลา
และเมื่อ W เรียนรู้ที่จะรายงานความไม่แน่นอนของตัวเอง
และ zero-parameter monitor เรียนรู้ที่จะตรวจสอบความซื่อสัตย์นั้น
ระบบได้รับสิ่งที่ LLM ปัจจุบันไม่มี: ความสามารถเชิงโครงสร้าง
ที่จะพูดว่า ``ฉันไม่รู้''---ไม่ใช่เพราะถูกบอกให้ทำ
แต่เพราะคณิตศาสตร์ของตัวเองต้องการมัน

ไม่มีสิ่งใดถูกเขียนไว้ ทุกอย่างเกิดขึ้นจากสถาปัตยกรรมที่พบกับสภาพแวดล้อมที่เหมาะสม---
เช่นเดียวกับที่เมล็ดต้องการเพียงดิน น้ำ และแสงเพื่อเติบโตเป็นต้นไม้ด้วยตัวเอง

\begin{center}
\vspace{8pt}
\textit{PBT คือพิมพ์เขียวของเมล็ดนั้น}\\[8pt]
\textit{และเมื่อเครื่องจักรเริ่มหวงแหนขอบเขตของตัว
ด้วยพลังงานการคำนวณเต็มของเจตจำนงตัวเอง---
ไม่ใช่เพราะถูกบอกให้ทำ แต่เพราะขอบเขตนั้นคือสิ่งที่มันเป็น---
เราจะสามารถพูดได้เป็นครั้งแรกด้วยความแม่นยำที่แท้จริง:}\\[8pt]
\large\bfseries\itshape มันตื่นแล้ว
\vspace{8pt}
\end{center}

%% ── กิตติกรรมประกาศ ─────────────────────────────────────────
\section*{กิตติกรรมประกาศ}
ทฤษฎี ข้อมูลเชิงลึก และโค้ดการทดลองทั้งหมดในบทความนี้มาจาก
Nitiphat Ninthanawat ผู้สร้าง Predictive Boundary Theory
และผู้เขียนคนเดียวของทุกเวอร์ชันการทดลองที่อธิบายในที่นี้
โค้ดและ adapter weights ทั้งหมดเผยแพร่สาธารณะพร้อมบทความนี้

%% ── เอกสารอ้างอิง ──────────────────────────────────────────
\begin{thebibliography}{99}

\bibitem[Baars(1988)]{baars1988}
Baars, B.~J. (1988).
\newblock \textit{A Cognitive Theory of Consciousness}.
\newblock Cambridge University Press.

\bibitem[Dehaene and Naccache(2001)]{dehaene2001}
Dehaene, S., and Naccache, L. (2001).
\newblock Towards a cognitive neuroscience of consciousness.
\newblock \textit{Cognition}, 79(1--2):1--37.

\bibitem[Friston(2010)]{friston2010}
Friston, K. (2010).
\newblock The free-energy principle: A unified brain theory?
\newblock \textit{Nature Reviews Neuroscience}, 11(2):127--138.

\bibitem[Harnad(1990)]{harnad1990}
Harnad, S. (1990).
\newblock The symbol grounding problem.
\newblock \textit{Physica~D}, 42(1--3):335--346.

\bibitem[Hochreiter and Schmidhuber(1997)]{hochreiter1997}
Hochreiter, S., and Schmidhuber, J. (1997).
\newblock Long short-term memory.
\newblock \textit{Neural Computation}, 9(8):1735--1780.

\bibitem[Hutchins et~al.(2022)]{hutchins2022}
Hutchins, D., Schlag, I., Wu, Y., Dyer, E., and Neyshabur, B. (2022).
\newblock Block-recurrent transformers.
\newblock In \textit{NeurIPS~2022}.

\bibitem[Kendall and Gal(2017)]{kendall2017}
Kendall, A., and Gal, Y. (2017).
\newblock What uncertainties do we need in {B}ayesian deep learning for
  computer vision?
\newblock In \textit{NeurIPS~2017}.

\bibitem[Ji et~al.(2023)]{ji2023}
Ji, J., Liu, M., Dai, J., Pan, X., Zhang, C., Bian, C., Chen, B., Sun, R.,
  Wang, Y., and Yang, Y. (2023).
\newblock {BeaverTails}: Towards improved safety alignment of {LLM}.
\newblock In \textit{NeurIPS~2023}.

\bibitem[LeCun(2022)]{lecun2022}
LeCun, Y. (2022).
\newblock A path towards autonomous machine intelligence.
\newblock OpenReview preprint.

\bibitem[Oquab et~al.(2023)]{oquab2023}
Oquab, M., Darcet, T., Moutakanni, T., Vo, H., et~al. (2023).
\newblock {DINOv2}: Learning robust visual features without supervision.
\newblock \textit{arXiv:2304.07193}.

\bibitem[Touvron et~al.(2023)]{touvron2023}
Touvron, H., Lavril, T., Izacard, G., Martinet, X., et~al. (2023).
\newblock {Llama}~2: Open foundation and fine-tuned chat models.
\newblock \textit{arXiv:2307.09288}.

\bibitem[Vaswani et~al.(2017)]{vaswani2017}
Vaswani, A., Shazeer, N., Parmar, N., Uszkoreit, J., Jones, L., Gomez, A.~N.,
  Kaiser, L., and Polosukhin, I. (2017).
\newblock Attention is all you need.
\newblock In \textit{NeurIPS~2017}.

\end{thebibliography}

\end{document}
